%%% Hlavní soubor. Zde se definují základní parametry a odkazuje se na ostatní části. %%%

%% Verze pro jednostranný tisk:
% Okraje: levý 40mm, pravý 25mm, horní a dolní 25mm
% (ale pozor, LaTeX si sám přidává 1in)
\documentclass[12pt,a4paper,draft]{report}
\setlength\textwidth{145mm}
\setlength\textheight{247mm}
\setlength\oddsidemargin{15mm}
\setlength\evensidemargin{15mm}
\setlength\topmargin{0mm}
\setlength\headsep{0mm}
\setlength\headheight{0mm}
% \openright zařídí, aby následující text začínal na pravé straně knihy
\let\openright=\clearpage

%% Pokud tiskneme oboustranně:
% \documentclass[12pt,a4paper,twoside,openright]{report}
% \setlength\textwidth{145mm}
% \setlength\textheight{247mm}
% \setlength\oddsidemargin{15mm}
% \setlength\evensidemargin{0mm}
% \setlength\topmargin{0mm}
% \setlength\headsep{0mm}
% \setlength\headheight{0mm}
% \let\openright=\cleardoublepage

%% Pokud používáte csLaTeX (doporučeno):
%\usepackage{czech}
% Pokud nikoliv:
\usepackage[czech]{babel}
\usepackage[T1]{fontenc}

%% Použité kódování znaků: obvykle latin2, cp1250 nebo utf8:
\usepackage[utf8]{inputenc}

%% Ostatní balíčky
\usepackage{tikz}
\usepackage{todonotes}
\usepackage{graphicx}
\usepackage{amsthm}
\usepackage{amsmath}
\usepackage{amsfonts}
\usepackage{epstopdf}
\usepackage{bm}
\usepackage{algpseudocode}
\usepackage{algorithm}
\usepackage{cite}
\usepackage{setspace}
\usetikzlibrary{bayesnet}

%% Balíček hyperref, kterým jdou vyrábět klikací odkazy v PDF,
%% ale hlavně ho používáme k uložení metadat do PDF (včetně obsahu).
%% POZOR, nezapomeňte vyplnit jméno práce a autora.
\usepackage[unicode]{hyperref}   % Musí být za všemi ostatními balíčky
\hypersetup{pdftitle=Název práce}
\hypersetup{pdfauthor=David Marek}

%%% Drobné úpravy stylu

% Tato makra přesvědčují mírně ošklivým trikem LaTeX, aby hlavičky kapitol
% sázel příčetněji a nevynechával nad nimi spoustu místa. Směle ignorujte.
\makeatletter
\def\@makechapterhead#1{
  {\parindent \z@ \raggedright \normalfont
   \Huge\bfseries \thechapter. #1
   \par\nobreak
   \vskip 20\p@
}}
\def\@makeschapterhead#1{
  {\parindent \z@ \raggedright \normalfont
   \Huge\bfseries #1
   \par\nobreak
   \vskip 20\p@
}}
\makeatother

% Toto makro definuje kapitolu, která není očíslovaná, ale je uvedena v obsahu.
\def\chapwithtoc#1{
\chapter*{#1}
\addcontentsline{toc}{chapter}{#1}
}

\renewcommand{\vec}[1]{\bm{#1}}
\hyphenation{prav-dě-po-do-bno-sti}
\hyphenation{sdru-že-né}

\floatname{algorithm}{Algoritmus}
\algblockdefx[INIT]{Init}{EndInit}{\textbf{init}}{\textbf{end init}}

\newtheorem{definice}{Definice}
\newtheorem{theorem}{Věta}


\begin{document}

% Trochu volnější nastavení dělení slov, než je default.
\lefthyphenmin=2
\righthyphenmin=2

%%% Titulní strana práce

\pagestyle{empty}
\begin{center}

\large

Univerzita Karlova v Praze

\medskip

Matematicko-fyzikální fakulta

\vfill

{\bf\Large DIPLOMOVÁ PRÁCE}

\vfill

\centerline{\mbox{\includegraphics[width=60mm]{../img/logo.eps}}}

\vfill
\vspace{5mm}

{\LARGE David Marek}

\vspace{15mm}

% Název práce přesně podle zadání
{\LARGE\bfseries Implementace aproximativních bayesovských metod pro odhad stavu v dialogových systémech}

\vfill

% Název katedry nebo ústavu, kde byla práce oficiálně zadána
% (dle Organizační struktury MFF UK)
Ústav formální a aplikované lingvistiky

\vfill

\begin{tabular}{rl}

Vedoucí diplomové práce: & Ing. Mgr. Filip Jurčíček, Ph.D. \\
\noalign{\vspace{2mm}}
Studijní program: & program \\
\noalign{\vspace{2mm}}
Studijní obor: & obor \\
\end{tabular}

\vfill

% Zde doplňte rok
Praha 2013

\end{center}

\newpage

%%% Následuje vevázaný list -- kopie podepsaného "Zadání diplomové práce".
%%% Toto zadání NENÍ součástí elektronické verze práce, nescanovat.

%%% Na tomto místě mohou být napsána případná poděkování (vedoucímu práce,
%%% konzultantovi, tomu, kdo zapůjčil software, literaturu apod.)

\openright

\noindent
Poděkování.

\newpage

%%% Strana s čestným prohlášením k diplomové práci

\vglue 0pt plus 1fill

\noindent
Prohlašuji, že jsem tuto diplomovou práci vypracoval(a) samostatně a výhradně
s~použitím citovaných pramenů, literatury a dalších odborných zdrojů.

\medskip\noindent
Beru na~vědomí, že se na moji práci vztahují práva a povinnosti vyplývající
ze zákona č. 121/2000 Sb., autorského zákona v~platném znění, zejména skutečnost,
že Univerzita Karlova v Praze má právo na~uzavření licenční smlouvy o~užití této
práce jako školního díla podle §60 odst. 1 autorského zákona.

\vspace{10mm}

\hbox{\hbox to 0.5\hsize{%
V ........ dne ............
\hss}\hbox to 0.5\hsize{%
Podpis autora
\hss}}

\vspace{20mm}
\newpage

%%% Povinná informační strana diplomové práce

\vbox to 0.5\vsize{
\setlength\parindent{0mm}
\setlength\parskip{5mm}

Název práce:
Implementace aproximativních bayesovských metod pro odhad stavu v dialogových systémech
% přesně dle zadání

Autor:
David Marek

Katedra:  % Případně Ústav:
Ústav formální a aplikované lingvistiky
% dle Organizační struktury MFF UK

Vedoucí diplomové práce:
Ing. Mgr. Filip Jurčíček, Ph.D., Ústav formální a aplikované lingvistiky
% dle Organizační struktury MFF UK, případně plný název pracoviště mimo MFF UK

Abstrakt:
% abstrakt v rozsahu 80-200 slov; nejedná se však o opis zadání diplomové práce

Klíčová slova:
% 3 až 5 klíčových slov

\vss}\nobreak\vbox to 0.49\vsize{
\setlength\parindent{0mm}
\setlength\parskip{5mm}

Title:
% přesný překlad názvu práce v angličtině

Author:
David Marek

Department:
Institute of Formal and Applied Linguistics
% dle Organizační struktury MFF UK v angličtině

Supervisor:
Ing. Mgr. Filip Jurčíček, Ph.D., Institute of Formal and Applied Linguistics
% dle Organizační struktury MFF UK, případně plný název pracoviště
% mimo MFF UK v angličtině

Abstract:
% abstrakt v rozsahu 80-200 slov v angličtině; nejedná se však o překlad
% zadání diplomové práce

Keywords:
% 3 až 5 klíčových slov v angličtině

\vss}

\newpage

%%% Strana s automaticky generovaným obsahem diplomové práce. U matematických
%%% prací je přípustné, aby seznam tabulek a zkratek, existují-li, byl umístěn
%%% na začátku práce, místo na jejím konci.

\openright
\pagestyle{plain}
\setcounter{page}{1}
\tableofcontents

%%% Jednotlivé kapitoly práce jsou pro přehlednost uloženy v samostatných souborech
\onehalfspacing
% !TEX root = prace.tex
\chapter*{Úvod}
\addcontentsline{toc}{chapter}{Úvod}

Dialog je přirozený způsob dorozumívání a sdělování informací mezi lidmi.
Počítač, který by dokázal vést dialog s uživatelem, byl vždy snem nejen
příznivců vědecko-fantastické literatury.
Už pro první počítače vnikaly programy, které se snažily využívat přirozenou
řeč pro interakci s uživatelem.
Jedním z takových programů byl například Eliza, program, který předstíral, že
jej zajímá, co mu uživatel říká.
Fungoval na principu rozpoznání textu pomocí gramatiky a následné transformace
textu do promluv dle pravidel.
Avšak gramatiky a pravidlové systémy se ukázaly nedostačné pro praktické
aplikace a tak se vývoj přesunul do statistických metod.
S využitím statistickým metod a metod strojového učení bylo možné začít s
porozumíváním mluveného slova.
Přijetí bylo zpočátku chladné a veřejnost byla 

\section{Cíle}

Dialogové systémy často využívají pouze nejlepší hypotézu ze systému porozumění přirozené řeči.
Většina ovšem umí vytvořit seznam $n$ nejlepších hypotéz.
Tato práce si klade za cíl představit metodu pro inferenci dialogového stavu v dialogovém systému s využitím více hypotéz.
Představená metoda bude založena na reprezentaci dialogového stavu pomocí dynamických bayesovských sítí.

Bude představen algoritmus Loopy Belief Propagation, který bude implementován pro využití v reálných systémech pro odhad dialogového stavu.
Algoritmus bude otestován na datech z reálného dialogového systému Let's Go a porovnán s dalšími systémy, které se účastnily soutěže Dialog State Tracking Challenge 2013.
Důležitou části algoritmů pro inferenci je určení pořadí, v jakém má inference probíhat.
Práce bude obsahovat implementaci několika strategií pro inferenci, které budou umožňovat efektivní inferenci pro různé druhy bayesovských sítí (stromy, dynamické sítě, obecné grafy).

Nakonec se práce bude zabývat učením parametrů sítě a představí algoritmus Expectation Propagation, který umožňuje inferenci v bayesovských sítích se spojitými náhodnými proměnnými.
Algoritmus Expectation Propagation je pouze konceptem, do kterého je třeba zasadit výpočet parametrů pravděpodobnostních rozložení z konkrétního grafického modelu.
Cílem práce není vytvořit obecný systém pro inferenci v libovolném grafickém modelu.
Takový cíl není ani reálný, protože při většině reálných použití dochází k aproximacím a úprávám v závislosti na ontologii problému, aby bylo možné inferenci provádět v reálném čase.
Cílem této práce tedy je vytvořit základní funkcionalitu potřebnou pro inferenci s EP, s kterou bude možné jednoduše vytvořit funkční systém, impolementovaný s doménově specifickými informacemi.
Jako příklad bude ukázán systém pro učení parametrů pravděpodobnostního rozložení pro pozorování v generickém modelu pro reprezentaci dialogového stavu.

% !TEX root = prace.tex
\chapter{Teorie dialogových systémů}

\section{Dialogový systém}

Dialogový systém je počítačový systém, který umožňuje uživatelům používat
hlasových pokynů pro ovládání počítačů a získávání informací. Dialogové systémy
nejsou pouhým hlasovým ovládáním počítače. Obsahují strategii řízení dialogu a
navíc obsahují interní stav dialogu, díky kterému jsou schopny vést dialog s
uživatelem. Dialog pak umožňuje využít složitějších příkazů, pokládat dotazy
nad strukturovanými daty a upřesňovat dotaz.

Dialogové systémy stále mají spoustu problémů k překonání a pro praktické použití je třeba se uchýlit k několika předpokladům a zjednodušením.
Prvním zjednodušením je doménová specializace, v současné době není možné vytvořit dialogový systém, který by se dokázal s uživatelem bavit o libovolném tématu. 
Vždy je potřeba při vývoji dialogového systému vědět, jaké informace má systém poskytovat a o čem se může chtít uživatel bavit.
Další zjednodušení se týkají přímo dialogu. 
Předpokládá se, že dialog probíhá vždy mezi systémem a jedním uživatelem. Navíc se tito pravidelně střídají. Jedna obrátka dialogu je složená z jedné promluvy systému a jedné promluvy uživatele.

Příkladem dialogového systému může být systém pro nalezení spojení pomocí
městské dopravy. Příkladem příkazu od uživatele pak může být např. \uv{Chci jet z
Malostranského náměstí na Anděl}. Dialogový systém se nyní může rozhodnout,
zda-li mu zadané informace stačí pro nalezení spojení. V tomto případě systém
stále neví, kdy chce uživatel jet. Může předpokládat, že uživatel už na
zastávce stojí a tak tedy uživateli nalezne nejbližší spojení.

Důležitou vlastností dialogového systému je robustnost. Pokud budeme používat
dialogový systém v přirozeném prostředí, tak se musíme vyrovnat s tím, že často
nebude uživateli rozumět. Uživatel může chtít jet stejně jako v minulém
případě, ale nyní dialogový systém přeslechne zastávku Anděl. Pro nalezení
spojení už nemá dost informací, může si pořád domyslet, že uživatel chce jet
nyní, ale je mnohem více možností, kterým směrem se vůbec chce uživatel vydat.
Dialogový systém se proto uživatele zeptá na chybějící informace a zároveň může
i implicitně potvrdit, že rozuměl alespoň odchozí zastávku: \uv{Dobře, jedete ze
zastávky Malostranské náměstí, můžete říct na jakou zastávku chcete jet?}.

\subsection{Typy dialogových systémů}

Dialogové systémy můžeme dělit podle několika kritérií. Podle toho, která
strana dialog vede anebo třeba jakým způsobem funguje strategie dialogového
systému.

\subsubsection{Dělení podle iniciativy}

Nejjednodušší dialogové systémy fungují na stejném principu jako automatické
telefonní systémy. Systém je ten, kdo vede dialog, a uživatel pouze odpovídá na
dotazy nebo vybírá z nabídnutých možností. Tento systém je velmi jednoduchý na
implementaci, pokud nutíme uživatel, aby vždy pouze odpověděl na položený
dotaz, tak velmi omezíme jeho vyjadřovací možnosti a díky tomu bude mnohem
jednodušší rozpoznat co řekl. Například pokud se uživatele ptáme jestli musí
nalezené spojení obsahovat pouze bezbariérové dopravní prostředky, tak můžeme
očekávat pouze \uv{Ano} nebo \uv{Ne} a jejich synonyma. Tento přístup ovšem
nebude velmi populární mezi častými uživateli systému, kteří by rádi zrychlili
postup dialogem.

Další možností je tedy povolit uživateli sdělit informace i tehdy, kdy se ho
systém ptá na něco jiného. Systém se tedy stále ptá na informace, které mu
chybějí, ale může dialog začít obecnější otázkou: "Jaké spojení chcete nalézt?"
oproti "Z jaké zastávky chcete jet?". V prvním případě dává systém pokročilému
uživateli možnost popsat celou cestu a velmi rychle se tak dostat ke kýžené
informaci. Pro nové uživatele, ale může být takováto obecně položená otázka
matoucí a tak je většinou třeba přidat i příklad.

Posledním typem je takový, kde uživateli přenecháme iniciativu, příkladem může
být například systém \uv{How may I help you?} \cite{gorin1997may} od
společnosti AT\&T. Ten slouží pro klasifikaci požadavků od uživatele a jejich
přepojení k patřičnému operátorovi.

\subsubsection{Dělení podle strategie řízení}

Důležité rozdělení dialogových systémů je podle způsobu, jakým funguje jejich
strategie řízení. Systém většinou předpokládá, že se v dialogu střídá s
uživatelem. Vždy, když na systém přijde řada, tak se musí rozhodnout, zda-li už
má dost informací k zodpovězení dotazu uživatele, případně zda-li se musí
dotázat na další informace, nechat si nějakou informaci potvrdit od uživatele,
anebo nechat uživatele vybrat si z nabízených možností.

Systém může mít přístup do databáze s informacemi a musí se tedy umět
rozhodnout, zda-li je třeba v ní hledat. Pak musí samozřejmě na nalezená data
reagovat. Může se stát, že je informací, které odpovídají zadání uživatele
příliš velké množství a tak je třeba zpřísnit zadání, anebo naopak nebyla žádná
informace nalezena a tedy je potřeba uživatele o tomto informovat.

Strategii, podle které se systém bude v těchto případech řídit lze
implementovat několika způsoby. Přímým řešením je napsat ji ručně. Pak se
rozhodnutí systému skládá z procházení série podmínek až do nalezení podmínky
odpovídající aktuálnímu stavu a následné provedení akce specifikované touto
podmínkou. Tento způsob může fungovat pro jednoduché systémy a pokud je vhodně
napsaný, tak může fungovat velmi dobře. Dost často ale tento způsob selhává,
pokud je dialogový systém příliš složitý, když se stávají podmínky
neudržitelné, anebo v případě, kdy se nemůžeme spolehnout na rozpoznanou řeč.

Další možností jak vytvořit strategii řízení je použít učení s učitelem.
Předpokládáme, že máme databázi hovorů mezi uživatelem a živým operátorem, z
kterých se dialogový systém může naučit, jak v dialogu postupovat. Tento
přístup se zdá jako dobrý nápad, ovšem naučíme se pouze napodobit operátora,
což nemusí být ten nejlepší způsob jak vést dialog. Navíc je nepřeberné
množství směrů, kterým se může dialog vyvíjet a jakmile se v jedné akci
odchýlíme, tak dostáváme úplně jiný dialog a tedy je velmi těžké z pozorovaných
dialogů zobecňovat. Proto se tento způsob téměř nevyužívá, existují ovšem
pokusy jak jej skloubit s následující metodou učení \cite{levin2000stochastic}.

Nejpoužívanější metodou v současnosti jak učit strategii dialogového systému
je použít zpětnovazební učení \cite{singh1999reinforcement}, \cite{walker2011application}. 
V takovém případě nám stačí pouze zhodnocení výsledku dialogu, které můžeme získat požádáním uživatelů dialogového systému, aby na konci dialogu oznámkovali, jak byli s výsledkem spokojeni. 
Při učení se nesnažíme napodobit existujícího operátora, ale získat co nejlepší hodnocení od uživatelů.
Díky tomu se může dialogový systém naučit pracovat dokonce ještě lépe než operátor.

\section{Součásti dialogového systému}

Dialogový systém se skládá z několika částí, každá se specializuje na jiný úkol a dohromady tvoří spolupracující systém.
Na vstupu je vždy zvukový záznam, o převedení do textu se stará systém rozpoznávání řeči (ASR).
Z textu je potřeba získat sémantické informace pomocí systému porozumění mluvené řeči (SLU).
Nad sémanticky anotovanými informacemi už může pracovat dialogový manager (DM), který rozhoduje o další akci.
Výstupem dialogového manageru jsou informace, které se mají předat uživateli.
O jejich převedení do textu se stará systém generování přirozené řeči (NLG).
Do zvukového záznamu převede text syntetizér řeči (TTS).

\subsection{Systém rozpoznávání řeči (ASR)}

Systém rozpoznávání řeči slouží k převedení mluveného projevu do textové podoby. 
Po získání textové podoby je teprve možné se zabývat významem textu.
Aktuálně nejlepší systémy jsou založené na pravděpodobnostním modelu a využívají skryté Markovské modely (HMM). 
Nejznámější je systém HTK \cite{young2002htk}, dalším otevřeným systémem je např. Kaldi \cite{Povey_ASRU2011}.
Existuje i celá řada komerčního software od firem jako IBM nebo Nuance.

Úspěšnost systému rozpoznání řeči je závislá na obtížnosti úlohy a na počtu trénovacích dat, pocházejících ze stejné domény.
Pro obecnou doménu se problém stává mnohem těžší.
Příkladem obecného rozpoznávače je Google Voice Search.

Systém rozpoznávání řeči může produkovat více hypotéz pro jeden vstup.
Často existuje pro jeden zvukový záznam více možných slovních sekvencí, z kterých by mohl pocházet.
Jejich reprezentace může být seznamem možných hypotéz s věrohodností pro každou hypotézu.
Věrohodnosti jsou skóre přiřazené hypotézám, které určují jakou důvěru má systém rozpoznávání řeči ve správnost dané hypotézy.
V ideálním případě je věrohodnost ekvivalentní aposteriorní pravděpodobnosti sekvence slov, dáno vstupní zvuk.
Ovšem ne všechny rozpoznávače pracují na pravděpodobnostním principu a pak není možné od nich požadovat skutečné pravděpodobnosti.

Další možností jak reprezentovat výstup je použití konfůzní sítě \cite{bertoldi2005new}. 
Konfůzní síť je vážený orientovaný graf obsahující startovní a konečný vrchol a hrany má označené slovy. 
Každá cesta ze startovního do konečného vrcholu vede přes všechny ostatní vrcholy. 
Váhy hran jsou pravděpodobnosti slova přiřazeného dané hraně. 
Hrany mohou obsahovat i prázdné slovo $\epsilon$.
Pravděpodobnost sekvence slov je součinem vah po cestě ze startovního do konečného uzlu.
Výhodou konfůzní sítě je, že umožňuje v komprimované podobě uložit mnohem více hypotéz.

\subsection{Porozumění mluvené řeči (SLU)}

Systém SLU převádí text do dialogových aktů. 
Dialogový akt (DA) je reprezentace promluvy, skládá se z jedné nebo více položek dialogového aktu (DAI), které jsou spojené v konjunkci.
Každá DAI se skládá z typu, názvu slotu a jeho hodnoty. 
Příklad DA z dialogového systému pro hledání restaurací:

\begin{center}
{\tt hello()\&inform(food="chinese")}.
\end{center}

Zde se DA skládá ze dvou položek, první položka má pouze typ {\em hello}, značící pozdrav. 
Druhá položka má typ {\em inform}, tzn. uživatel nás informuje o svém požadavku.
Název slotu je {\em food} a hodnota je \uv{chinese}, tedy uživatel nám říká, že hledá restauraci, kde servírují čínské jídlo.

Typů může být libovolné množství, ale existuje několik základních, jejichž použití je ustálené. 
\begin{itemize}
\item {\em inform} --- sdělujeme informaci, doplňujeme hodnotu do slotu,
\item {\em request} --- požadujeme od protějšku doplnění hodnoty pro dotazovaný slot,
\item {\em confirm} --- chceme potvrdit hodnotu slotu, potvrzení může být implicitní, anebo explicitní. 
	Při explicitním potvrzení očekáváme odpověď buď \uv{Ano} nebo \uv{Ne},
	U implicitního, pokud se nám nedostane odpovědi předpokládáme, že protějšek souhlasí,
\item {\em select} --- žádáme protějšek, aby zvolil z nabízených možností.
\end{itemize}

\subsection{Dialogový manager (DM)}

Dialogový manager tvoří mozek dialogového systému, jeho vstupem jsou dialogové akty z SLU, výstupem opět dialogové akty, pro NLG.
DM se dělí na dvě části. 
První část se stará o udržování dialogového stavu, tzn. informace, které uživatel systému poskytl, historii dialogu, které informace už byly poskytnuty anebo např. potvrzeny či zamítnuty uživatelem.
Dále obsahuje dialogovou strategii, která určuje příští akci v závislosti na dialogovém stavu.

\subsubsection{Reprezentace stavu dialogu}

Markovský rozhodovací proces (MDP) \cite{puterman2009markov} je diskrétní, stochastický a kontrolovaný proces. 
V každém časovém okamžiku se systém nachází v nějakém stavu $s$. 
Uživatel provede akci $a$, dostupnou ve stavu $s$, a ta jej přesune náhodně do nového stavu $s^\prime$ a navíc dostane odměnu $r(s, s^\prime, a)$.
Pravděpodobnost přechodu ze stavu $s$ do stavu $s^\prime$ je dána přechodovou funkcí $p(s^\prime \mid s, a)$.

Zobecněním MDP je částečně pozorovatelný Markovský rozhodovací proces (POMDP). 
Částečně pozorovatelný proto, že na rozdíl od MDP nevíme v jakém stavu se systém nachází.
Naše představa o stavu je dána pouze pravděpodobnostním rozložením přes všechny možné stavy, tzv. belief $b(s)$.
Výsledek provedení akce tedy nezávisí pouze na přechodové pravděpodobnosti, ale také na belief stavu. 
Pro aktualizaci stavu je potřeba vysčítat přes všechny možné stavy
\begin{equation}
b(s^\prime) = \sum_s p(s^\prime \mid s, a) b(s).
\end{equation}

Dialog je tedy popsán pomocí POMDP \cite{williams2007partially}, protože musíme zakomponovat naši neznalost cílů uživatele a problémy s rozpoznáváním. 
Pro reprezentaci stavu dialogu existuje několik různých přístupů.

\section{Grafické modely}


% !TEX root = prace.tex
\chapter{Bayesovské sítě a inference}
\label{ch:kap2}

V této kapitole představíme Bayesovské sítě, grafický model pro efektivní reprezentaci pravděpodobnostních rozdělení a nezávislostí mezi náhodnými proměnnými.
Bayesovská síť zároveň vytváří koncept pro inferenci, tedy zodpovídání dotazů nad proměnnými v síti.
Ukážeme si několik přístupů k inferenci, nejprve naivní výpočet vycházející přímo z definice.
Následně využijeme vlastností sítě a konceptů dynamického programování pro jeho zlepšení.
Analýzou složitosti exaktní inference dojdeme k závěru, že pro větší a komplexnější modely bude třeba se uchýlit k aproximacím.

Nejprve aproximujeme sdruženou pravděpodobnostní distribuci součinem marginálních distribucí a představíme Loopy Belief Propagation algoritmus.
Stále můžeme použít pouze diskrétní pravděpodobnostní distribuce se známými parametry.
Pro učení parametrů lze použít Expectation Maximization metodu, pro dialogové systémy je ovšem těžké získat dostatek učících dat.
Dostaneme se k metodě Expectation Propagation, která je zobecněním LBP a je možné ji použít pro libovolné rozdělení.
Díky ní budeme schopní vytvořit generativní model pro aktualizaci dialogového stavu, který bude pracovat stejně jako LBP, ale bude schopen adaptace za běhu dialogového systému.

\section{Bayesovské sítě}
\label{sec:bn}

Bayesovské sítě jsou pravděpodobnostní grafický model, který využívá podmíněných nezávislostí pro úspornou reprezentaci sdružené pravděpodobnosti.
Bayesovská sít je orientovaný acyklický graf, jeho vrcholy jsou náhodné proměnné a hrany odpovídají přímé závislosti jednoho uzlu na druhý.
Pro každou náhodnou proměnnou v síti platí, že její pravděpodobnost je jednoznačně určena jejími rodiči v grafu.
Podmíněná pravděpodobnostní distribuce (CPD) proměnné $X$ popisuje pravděpodobnost proměnné $X$ dáno její rodiče, $P(X \mid parents(X))$.
Pokud proměnná nemá žádná rodiče, pak její podmíněná pravděpodobnostní distribuce je ekvivalentní marginální pravděpodobnostní distribuci.

Příklad Student~\cite{koller2009probabilistic}: firma zvažuje, zda-li přijme studenta.
Firma chce přijímat chytré studenty, ale nesmí je testovat na inteligenci (I) přímo.
Má však výsledek studentových SAT testů, které ovšem nemusí stačit pro správné zhodnocení inteligence.
Požadují tak tedy i doporučení (D) od jednoho z učitelů.
Učitel studentovi napíše doporučující dopis na základě známky (Z), kterou student získal v jeho předmětu.
Předměty se ovšem liší v obtížnosti (O) a tak je studentova známka v předmětu závislá nejen na jeho inteligenci, ale také na obtížnosti předmětu.
Grafický model reprezentující tento problém je vyobrazen na obrázku~\ref{fig:student}.
\begin{figure}
\begin{center}
\begin{tikzpicture}

\matrix[row sep=0.75cm, column sep=1.2cm]
{
    \node[latent]       (O)     {O};
    && \node[latent]    (I)     {I};
    \\
    &\node[latent]      (Z)     {Z};
    &&\node[latent]     (S)     {SAT};
    \\
    &\node[latent]      (D)     {D};
    \\
};

\edge{O}{Z}
\edge{I}{Z}
\edge{I}{S}
\edge{Z}{D}

\end{tikzpicture}
\end{center}
\label{fig:student}
\caption{Bayesovská síť pro příklad se studentem.}
\end{figure}

V tomto modelu je několik nezávislostí. Obtížnost předmětu a inteligence studenta jsou zjevně nezávislé.
Studentova známka z předmětu je závislá na obtížnosti předmětu a inteligenci studenta, ale je podmíněně nezávislá na jeho výsledku ze SAT, dáno studentova inteligence.
Konečně doporučení, které student obdrží, je podmíněně nezávislé na všech ostatních proměnných, dáno studentova známka.

Sdruženou nezávislost tohoto modelu lze zapsat ve formě podmíněných pravděpodobnostních distribucí s pomocí řetízkového pravidla.
\begin{equation}
P(O, I, Z, S, D) = P(D \mid Z) P(Z \mid O, I) P(SAT \mid I) P(O) P(I)
\end{equation}

Předpokládejme, že obtížnost předmětu, inteligence studenta, doporučující dopis a výsledek SAT jsou binární proměnné.
Známka z předmětu pak je ternární proměnná.
Pokud bychom zapsali sdruženou pravděpodobnost ve formě tabulky, tak se dostaneme k 48 položkám.
Díky rozdělení do podmíněných pravděpodobnostních rozložení, které nám bayesovská síť poskytuje, se dostáváme k $2 + 2 + 12 + 4 + 6 = 26$ položkám.
Tedy i v tomto jednoduchém modelu dochází k značné úspoře.

\section{Inference v Bayesovských sítích}

Bayesovské sítě reprezentují pravděpodobnostní model a umožňují nám nad ním provádět dotazy.
Můžeme se například ptát na marginální pravděpodobnost jednotlivých proměnných. 
Tu získáme marginalizací sdružené pravděpodobnosti, pokud vezmeme příklad se studentem a budeme chtít znát marginální pravděpodobnost známek, musíme vysčítat všechny ostatní proměnné.

\begin{equation}
P(Z) = \sum_{O, I, S, D} P(D \mid Z) P(Z \mid O, I) P(SAT \mid I) P(O) P(I)
\end{equation}

Další a asi nejčastější dotaz nastává, pokud některé náhodné proměnné pozorujeme. 
Pak chceme znát pravděpodobnost jiných proměnných dáno naše pozorování, $P(\vec{X} \mid \vec{E} = \vec{e})$, kde $\vec{X}$ jsou dotazované proměnné, $\vec{E}$ jsou pozorované proměnné a $\vec{e}$ jsou pozorované hodnoty.
Z definice podmíněné pravděpodobnosti dostáváme

\begin{equation}
P(\vec{X} \mid \vec{E} = \vec{e}) = \frac{P(\vec{X}, \vec{E} = \vec{e})}{P(\vec{E} = \vec{e})}
\end{equation}

Každou instanci jmenovatele $P(\vec{X}, \vec{E} = \vec{e})$ jde spočítat sumou sdružených pravděpodobností s ohodnoceními proměnných, které jsou kompatibilní s pozorováním a aktuální instancí. 
Pokud počítáme instanci $P(\vec{X} = \vec{x}, \vec{E} = \vec{e})$, pak získáme výsledek marginalizací všech proměnných, kromě $\vec{X}$ a $\vec{E}$, které jsou fixovány.
Pokud tedy množinu všech proměnných bez $\vec{X}$ a $\vec{E}$ označíme $\mathcal{W} = \mathcal{X} - \vec{X} - \vec{E}$, pak pravděpodobnost dané instance je

\begin{equation}
P(\vec{X} = \vec{x}, \vec{E} = \vec{e}) = \sum_{\vec{w}} P(\vec{x}, \vec{e}, \vec{w})
\end{equation}

Pro výpočet normalizační konstanty $P(\vec{E})$ musíme opět marginalizovat sdruženou pravděpodobnost, anebo si můžeme povšimnout, že platí
\begin{equation}
P(\vec{E} = \vec{e}) = \sum_{\vec{x}} P(\vec{x}, \vec{e})
\end{equation}
a tedy můžeme použít už vypočítané hodnoty.

\subsection{Exaktní inference}

V předchozí části jsme viděli, že pomocí definice podmíněné pravděpodobnosti a marginalizace sdružené pravděpodobnosti lze najít odpověď na libovolný dotaz. 
Nyní si ukážeme algoritmus, který využívá struktury Bayesovské sítě pro inferenci a navíc díky metodám dynamického programování umožňuje samotný výpočet urychlit.
Nakonec ovšem zjistíme, že pro velké sítě, které nás většinou zajímají nejvíce, nám přesná inference nebude stačit a musíme se uchýlit k aproximacím.

Začneme s inferencí v jednoduchém modelu $A \rightarrow B \rightarrow C \rightarrow D$.
Sdružená pravděpodobnost $P(A, B, C, D)$ je součinem jednotlivých podmíněných pravděpodobnostních distribucí
\begin{equation}
P(A, B, C, D) = P(D \mid C) P(C \mid B) P(B \mid A) P(A)
\end{equation}

Pokud nyní budeme chtít spočítat marginální distribuci $D$, tak musíme marginalizovat všechny ostatní proměnné
\begin{equation}
P(D) = \sum_{A, B, C} P(D \mid C) P(C \mid B) P(B \mid A) P(A)
\end{equation}

Můžeme si povšimnout, že spousta členů se bude počítat vícekrát.
Využitím metod dynamického programování a přeuspořádáním sum si můžeme mezivýsledky uložit a použít vícekrát.
\begin{equation}
P(D) = \sum_C P(D \mid C) \sum_B P(C \mid B) \sum_A P(B \mid A) P(A)
\end{equation}

Při výpočtu pak nejprve spočítáme $\psi_1(A, B) =  P(B \mid A) P(A)$,  pak vysčítáme proměnnou $A$ a získáme $\tau_1(B) = \sum_A \psi_1(A, B)$.
Pokračujeme obdobně
\begin{align}
    \psi_2(B, C) &= P(C \mid B) \tau_1(B) \\
    \tau_2(C) &= \sum_B \psi_2(B, C)
\end{align}
A nakonec spočítáme finální marginální pravděpodobnost
\begin{align}
    \psi_3(D, C) &= P(D \mid C) \tau_2(C) \\
    P(D) &= \sum_C \psi_3(D, C)
\end{align}

Nechť $\mathcal{X}$ je množina náhodných proměnných. 
$Val(\mathcal{X})$ definujeme jako množinu možných ohodnocení proměnných $\mathcal{X}$.
Potom definujeme faktor $\phi$ jako zobrazení z $Val(\mathcal{X})$ do $\mathbb{R}$. Faktor je nezáporný, pokud všechny jeho obrazy jsou nezáporné. Množina proměnných $\mathcal{X}$ je doménou faktoru a značíme ji jako $Dom(\phi)$.

Faktor, jehož doménu tvoří diskrétní proměnné, si můžeme představit jako tabulku, která obsahuje jednu hodnotu pro každé možné ohodnocení proměnných z domény.

Nechť $\vec{X}, \vec{Y}, \vec{Z}$ jsou tři disjunktní množiny náhodných proměnných.
Nechť $\phi_1(\vec{X}, \vec{Y})$ a $\phi_2(\vec{Y}, \vec{Z})$ jsou faktory.
Definujeme součin faktorů $\phi_1 \times \phi_2$ jako faktor $\psi: Val(\vec{X}, \vec{Y}, \vec{Z}) \rightarrow \mathbb{R}$ následovně:
\begin{equation*}
\psi(\vec{X}, \vec{Y}, \vec{Z}) = \phi_1(\vec{X}, \vec{Y}) \cdot \phi_2(\vec{Y}, \vec{Z})
\end{equation*}

Násobíme prvky, které mají stejné ohodnocení společných proměnných $\vec{Y}$.
Stejný princip použijeme pro všechny matematické operace.

Nechť $\vec{X}$ je množina náhodných proměnných a $Y \not\in \vec{X}$ náhodná proměnná. 
Nechť $\phi(\vec{X}, Y)$ je faktor.
Definujeme marginalizaci $Y$ v $\phi$, značenou $\sum_Y \phi$, jako faktor $\psi$ s doménou $\vec{X}$ takový, že
\begin{equation*}
    \psi(\vec{X}) = \sum_Y \phi(\vec{X}, Y)
\end{equation*}
Této operaci také říkáme vysčítání $Y$ ve $\phi$.

Faktory se při počítání chovají jako čísla, všechny operace probíhají po prvcích, pro které je ohodnocení náhodných proměnných z průniku domén faktorů stejné.
Proto platí komutativita $\phi_1 \cdot \phi_2 = \phi_2 \cdot \phi_1$ a $\sum_X \sum_Y \phi = \sum_Y \sum_X \phi$.
Dále platí asociativita součinu $(\phi_1 \cdot \phi_2) \cdot \phi_3 = \phi_1 \cdot (\phi_2 \cdot \phi_3)$.
Nakonec můžeme vyměnit sumu a součin, pokud $X \not\in Dom(\phi_1)$, potom $\sum_X (\phi_1 \cdot \phi_2) = \phi_1 \sum_X \phi_2$.

Sdruženou pravděpodobnost z minulého příkladu tedy můžeme přepsat do formy faktorů.
\begin{equation}
    P(A, B, C, D) = \phi_A \cdot \phi_B \cdot \phi_C \cdot \phi_D
\end{equation}

Opět se pokusíme spočítat marginální pravděpodobnost proměnné $D$.
\begin{align}
P(D) &= \sum_C \sum_B \sum_A \phi_A \phi_B \phi_C
\\
&= \sum_C 
	\phi_D \cdot 
	\left( 
		\sum_B 
			\phi_C \cdot 
			\left( 
				\sum_A 
					\phi_A \cdot \phi_B
			\right)
	\right)
\end{align}

Přesuny sum můžeme provést díky doméně jednotlivých faktorů.
Faktory $\phi_C$ a $\phi_D$ neobsahují proměnnou $A$ a tedy je můžeme vytknout před sumu přes $A$.
Stejně tak faktor $\phi_D$ neobsahuje proměnnou $B$ a opět jej můžeme vytknout před sumu přes $B$.
Tyto úpravy můžeme provádět v libovolném pořadí, pokud vždy platí, že vysčítáme proměnnou $X$ až poté, co spolu vynásobíme všechny faktory, které ji obsahují.

V obecnosti vždy počítáme výraz, který je ve tvaru
\[
\sum_X \prod_{\phi \in \Phi} \phi.
\]

Z tohoto také vychází název pro tuto metodu: sum-product.
Jednoduchý algoritmus pro exaktní inferenci využívající tuto metodu se nazývá eliminace proměnných.
Základní myšlenka je, že máme dán seznam náhodných proměnných v pořadí, v jakém se mají eliminovat.
Pro eliminaci proměnné je třeba nejprve vynásobit všechny faktory, které ji obsahují a následně ji vysčítat.
Tak získáme faktor, který už tuto proměnnou neobsahuje, tedy jsme ji eliminovali.
Eliminace proměnných je popsána v algoritmu~\ref{alg:ve}.

\begin{algorithm}[H]
\caption{Eliminace proměnných}
\label{alg:ve}
\begin{algorithmic}
\Function{Sum-Product-VE}{$\Phi$, $\vec{X}$, $\prec$}
\State{$\Phi$ množina všech faktorů.} 
\State{$\vec{X}$ množina náhodných proměnných, které mají být eliminovány.}
\State{$\prec$ pořadí proměnných, v jakém mají být eliminovány.}
\State

\State Nechť $X_1, \dots, X_k$ je seřazení proměnných z $\vec{X}$, t.ž. $X_i \prec X_j \Leftrightarrow i < j$.
\For{$i = 1 \dots k$}
	\State $\Phi \gets$ Sum-Product-Eliminate-Var($\Phi$, $Z_i$)
\EndFor
\State $\phi^* \gets \prod_{\phi \in \Phi} \phi$
\State \Return $\phi^*$
\EndFunction
\State
\Function{Sum-Product-Eliminate-Var}{$\Phi$, $X$}
\State{$\Phi$ množina všech faktorů.}
\State{$X$, proměnná, která má být eliminována.}
\State

\State $\Phi^\prime \gets \{\phi \in \Phi: X \in Dom(\phi)\}$
\State $\Phi^{\prime \prime} \gets \Phi - \Phi^\prime$
\State $\psi \gets \prod_{\phi \in \Phi^\prime} \phi$
\State $\tau \gets \sum_X \psi$
\State \Return $\Phi^{\prime\prime} \bigcup \{\tau\}$
\EndFunction
\end{algorithmic}
\end{algorithm}

Nechť $\vec{X}$ je množina náhodných proměnných, nechť $\Phi$ je množina faktorů, t.ž. pro každé $\phi \in \Phi$, $Dom(\phi) \subseteq \vec{X}$.
Nechť $\vec{Y} \subset \vec{X}$ je množina dotazovaných náhodných proměnných a nechť $\vec{Z} = \vec{X} - \vec{Y}$.
Pak pro každé seřazení $\prec$ nad $\vec{Z}$, Sum-Product-VE($\Phi$, $\vec{Z}$, $\prec$) vrátí faktor $\phi^*(\vec{Y})$ takový, že
\[
\phi^*(\vec{Y}) = \sum_Z \prod_{\phi \in \Phi} \phi
\]

Nyní provedeme analýzu algoritmu eliminace proměnných.
Předpokládejme, že na vstupu je $n$ proměnných.
Bayesovská síť obsahuje pro každou proměnnou jeden faktor.
Pro jednoduchost budeme předpokládat, že algoritmus bude eliminovat všechny proměnné.
Běh algoritmu se skládá z jednotlivých eliminačních kroků, při kterých je vždy eliminována jedna proměnná.

Při jednom eliminačním kroku je vybrána proměnná $X_i$, všechny faktory, které ji obsahují jsou vynásobeny a vytvoří jeden velký faktor $\psi_i$.
Proměnná $X_i$ je pak vysčítána z tohoto faktoru.
Počet operací pro jeden eliminační krok tedy závisí na velikosti faktoru $\psi_i$, označme ji $N_i$.
Maximum z velikostí faktorů označme $N_{max} = max_i N_i$.

Nyní se zaměříme na počet násobení.
Celkem vznikne $n+m$ faktorů, kde $m$ je počet faktorů, které vznikly vysčítáním proměnné.
Každý z těchto faktorů je zahrnut do součinu pouze jednou, pokud je eliminována nějaká proměnná, kterou obsahuje.
Cena násobení faktorů pro vznik $\psi_i$ je nejvýše $N_i$.
Celkový počet násobení tedy bude nejvýše $(n+m)N_{max}$ což je $\mathcal{O}(nN_{max})$.

Pokud $k$ je maximum z velikostí domén proměnných, pak velikost faktoru obsahujícího $n$ proměnných může být až $k^n$.
Složitost eliminace proměnných je tedy dominována velikostí faktorů, které vznikají při výpočtu a je exponenciální.
Navíc bylo dokázáno, že výběr nejlepšího pořadí proměnných pro eliminaci je NP těžký~\cite{arnborg1987complexity}.

\subsection{Posílání zpráv ve faktor grafu}

Algoritmy exaktní inference naráží při reálném použití na příliš velkou složitost násobení faktorů.
Hlavním problémem je velikost sdružené pravděpodobnosti, která roste exponenciálně s počtem náhodných proměnných.
Většinou nás ovšem zajímá marginální pravděpodobnost jedné nebo jen mála proměnných.
Je tedy zbytečné počítat celou sdruženou pravděpodobnost, abychom z ní pak vysčítali většinu proměnných.
Řešení se nabízí ve formě aproximace faktorizací sdružené pravděpodobnosti.

\begin{equation*}
    P(X_1, \dots, X_n) = \prod_i P(X_i)
\end{equation*}

Pro výpočet s faktorizovanou distribucí si zavedeme novou datovou strukturu, tzv. faktor graf.
Faktor graf je bipartitní graf, kde jednu partitu tvoří faktory a druhou partitu tvoří náhodné proměnné.
Hrany ve faktor grafu vedou vždy mezi proměnnou a faktorem, který ji obsahuje.

\begin{figure}[H]
\begin{center}
\begin{tikzpicture}
    
\matrix[row sep=0.75cm, column sep=1.2cm]
{
    \node[latent]   (X1)    {$X_1$};
    &
    \factor         {fa}    {$f_a$}{}{};
    &
    \node[latent]   (X2)    {$X_2$};
    &
    \factor         {fb}    {$f_b$}{}{};
    &
    \node[latent]   (X3)    {$X_3$};
    \\
    &&
    \factor         {fc}    {left:$f_c$}{}{};
    \\
    &&
    \node[latent]   (X4)    {$X_4$}{}{};
    \\
};
\edge[-]{X1,X2} {fa}
\edge[-]{X3,X2} {fb}
\edge[-]{X2,X4} {fc}

\end{tikzpicture}
\end{center}
\caption{Příklad faktor grafu se třemi náhodnými proměnnými $X_1, X_2, X_3$ a třemi faktory $f_a, f_b, f_c$.}
\end{figure}

Nejprve začneme s inferencí na stromech.
Pokud zafixujeme jednu náhodnou proměnnou $X$ ve faktor grafu, pak sdruženou pravděpodobnost můžeme spočítat jako
\begin{equation}
    P(\vec{X}) = \prod_{s \in ne(X)} F_s(X, \vec{X}_s),
\label{eq:fs}
\end{equation}
kde $ne(X)$ jsou faktory obsahující proměnnou $X$ (tedy sousedi ve faktor grafu), $F_s$ je součin všech faktorů v podstromu určeném faktorem $f_s$ a $\vec{X}_s$ je množina všech proměnných v daném podstromu.

Pro výpočet marginální pravděpodobnosti $X$ substituujeme~(\ref{eq:fs}) do výpočtu marginální pravděpodobnosti ze sdružené a po výměně sumy a produktu dostaneme
\begin{align}
P(X) &= \prod_{s \in ne(X_i)} \sum_{\vec{X}_s} F_s(X, \vec{X}_s)
\\
&= \prod_{s \in ne(X)} \mu_{f_s \rightarrow X}(X)
\label{eq:margx}
\end{align}

Zavedli jsme funkce
\begin{equation}
    \mu_{f_s \rightarrow X} \equiv \sum_{\vec{X}_s} F_s(X, \vec{X}_s),
\label{eq:defmsgfx}
\end{equation}
které můžeme nazývat zprávami z faktoru $f_s$ do proměnné $X$.

Každé $F_s(X, \vec{X}_s)$ je popsáno podgrafem faktor grafu a tedy může být znova faktorizováno.
Namísto proměnné $X$ nyní vezmeme faktor $f_s$.
Náhodné proměnné sousedící s faktorem $f_s$ bez $X$ si označíme $X_1, \dots, X_M$.
Součin faktorů v podstromech určených těmito proměnnými označíme $G_i(X_i, \vec{X}_{si})$.
Faktor $F_s(X, \vec{X}_s)$ tedy můžeme přepsat jako
\begin{equation}
F_s(X, \vec{X}_s) = f_s(X, X_1, \dots, X_M) G_1(X_1, \vec{X}_{s1}), \dots, G_M(X_M, \vec{X}_{sM})
\end{equation}

Pokud substitujeme přepsaný faktor $F_s(X, \vec{X}_s)$ do definice zprávy z faktoru, dostaneme
\begin{align}
\mu_{f_s \rightarrow X}(X) &=
    \sum_{X_1} \dots \sum_{X_M}
        f_s(X, X_1, \dots, X_M)
        \prod_{m \in ne(f_s) \backslash X}
            \sum_{\vec{X}_{sm}}
                G_m(X_m, \vec{X}_{sm})
\\
&= \sum_{X_1} \dots \sum_{X_M}
    f_s(X, X_1, \dots, X_M)
    \prod_{m \in ne(f_s) \backslash X}
        \mu_{X_m \rightarrow f_s}(X_m)
\label{eq:mfsx}
\end{align}

Zavedli jsme další funkce 
\begin{equation}
    \mu_{X_m \rightarrow f_s}(X_m) \equiv \sum_{\vec{X}_{sm}} G_m(X_m, \vec{X}_{sm}),
\label{eq:defmsgxf}
\end{equation}
které budeme nazývat zprávami z proměnné $X_m$ do faktoru $f_s$.

Z rovnice~\ref{eq:margx} vidíme, že marginální pravděpodobnost proměnné vypočítáme jako součin zpráv ze všech okolních faktorů.
Každou z těchto zpráv můžeme spočítat jako součin faktoru a zpráv přicházejících z proměnných, které s tímto faktorem sousedí, kromě proměnné, které chceme zprávu posílat. Zbývá nám tedy zjistit v jakém tvaru jsou zprávy z proměnných do faktoru.

Stejně jako u $F_s(X, \vec{X}_s)$ nám tady $G_m(X_m, \vec{X}_{sm})$ definuje podgraf faktor grafu. 
V kořeni tohoto podgrafu leží proměnná $X_m$ a tím se dostáváme na už známý případ.
Budeme ignorovat faktor $f_s$, protože neleží v podgrafu určeném $G_m(X_m, \vec{X}_{sm})$, a pak 
\begin{equation}
    G_m(X_m, \vec{X}_{sm}) = \prod_{l \in ne(X_m) \backslash f_s} F_l(X_m, \vec{X}_{ml})
\end{equation}
Přepsáno ve formě zpráv
\begin{align}
\mu_{X_m \rightarrow f_s}(X_m)
&= \prod_{l \in ne(X_m) \backslash f_s}
    \sum_{\vec{X}_{ml}}
        F_l(X_m, \vec{X}_{ml})
\\
&= \prod_{l \in ne(X_m) \backslash f_s}
    \mu_{f_l \rightarrow X_m}(X_m)
\end{align}
Zpráva z náhodné proměnné je tedy součinem zpráv ze všech ostatních faktorů.

Zprávy tedy dokážeme počítat rekurzivně, chybí nám ovšem pravidla pro zprávy z uzlů, které jsou listy.
V případě faktoru je odchozí zpráva ekvivalentní faktoru.
\begin{equation}
    \mu_{f \rightarrow X}(X) = f(X)
\end{equation}
V případě náhodné proměnné je odchozí zpráva 
\begin{equation}
    \mu_{X \rightarrow f}(X) = 1
\end{equation}
Pro pozorovanou proměnnou je zpráva vždy rovna pozorované hodnotě, bez ohledu na sousední faktory.

\subsection{Belief Propagation}

Algoritmus exaktní inference na stromech s pomocí posílání zpráv se nazývá belief propagation~\cite{pearl1988probabilistic}.
Zpráva z vrcholu faktor grafu může být poslána do jiného vrcholu, až když byly obdrženy zprávy ze všech ostatních vrcholů.
Jakmile vrchol obdržel zprávy ze všech vrcholů, lze spočítat jeho marginální pravděpodobnost.
Strom můžeme zakořenit v libovolné náhodné proměnné a posílat zprávy z listů do kořene.
Kořen pak obdrží všechny zprávy ze všech sousedních faktorů a tedy je možné spočítat jeho marginální pravděpodobnost.
Tomuto se říká dopředný krok propagace.

\begin{figure}[H]
\begin{center}
\begin{tikzpicture}

\matrix[row sep=0.75cm, column sep=1.2cm]
{
    \node[latent]   (X1)    {$X_1$};
    &
    \factor         {fa}    {$f_a$}{}{};
    &
    \node[latent]   (X2)    {$X_2$};
    &
    \factor         {fb}    {$f_b$}{}{};
    &
    \node[latent]   (X3)    {$X_3$};
    \\
    &&
    \factor         {fc}    {left:$f_c$}{}{};
    \\
    &&
    \node[latent]   (X4)    {$X_4$}{}{};
    \\
};
\edge{X1}{fa}
\edge{fa}{X2}
\edge{X4}{fc}
\edge{fc}{X2}
\edge{X2}{fb}
\edge{fb}{X3}

\end{tikzpicture}
\end{center}
\caption{Směr posílání zpráv pro strom zakořeněný ve vrcholu $X_3$ v dopředném kroce.}
\end{figure}

Belief propagation algoritmus tedy umožňuje na stromech spočítat marginální pravděpodobnost jedné proměnné s linárním počtem poslaných zpráv.
Při rozšíření na spočítání marginální pravděpodobnosti všech proměnných není třeba algoritmus pouštět $n$-krát.
Stačí si uvědomit, že všem vrcholům chybí pouze zpráva od faktoru na cestě ke kořenu.
A my můžeme z kořene poslat zprávy zpět k listům hned, jak obdržíme všechny příchozí zprávy.
Této části se říká zpětný krok propagace.

\begin{figure}[H]
\begin{center}
\begin{tikzpicture}

\matrix[row sep=0.75cm, column sep=1.2cm]
{
    \node[latent]   (X1)    {$X_1$};
    &
    \factor         {fa}    {$f_a$}{}{};
    &
    \node[latent]   (X2)    {$X_2$};
    &
    \factor         {fb}    {$f_b$}{}{};
    &
    \node[latent]   (X3)    {$X_3$};
    \\
    &&
    \factor         {fc}    {left:$f_c$}{}{};
    \\
    &&
    \node[latent]   (X4)    {$X_4$}{}{};
    \\
};
\edge{fa}{X1}
\edge{X2}{fa}
\edge{fc}{X4}
\edge{X2}{fc}
\edge{fb}{X2}
\edge{X3}{fb}

\end{tikzpicture}
\end{center}
\caption{Směr posílání zpráv pro strom zakořeněný ve vrcholu $X_3$ ve zpětném kroce.}
\end{figure}

Po propagaci všech zpráv až k listům už každá proměnná získala zprávy od všech sousedních faktorů a tedy je možné spočítat marginální pravděpodobnost všech proměnných.

\subsection{Loopy Belief Propagation}
\label{sec:lbp}

Algoritmus Belief Propagation funguje pro inferenci na stromech.
Hlavní důvod proč nelze použít Belief Propagation na obecných grafech je, že můžeme narazit na cyklus v grafu, což znamená, že žádný z vrcholů v tomto cyklu nebude nikdy mít dostatek příchozích zpráv, aby mohl nějakou zprávu odeslat.

Algoritmus Loopy Belief Propagation řeší tento problém relaxací podmínky na příchozí zprávy.
Pro odeslání zprávy do uzlu není třeba znát zprávy ze všech ostatních sousedních vrcholů.
Chybějící zprávy jsou nastaveny na jedničku.
Loopy Belief Propagation je iterativní algoritmus.
V grafu nemusí existovat vrchol, který může posílat zprávu k nějakému ze svých sousedů.
Je tedy třeba vybrat vrchol podle nějaké strategie a z něj poslat zprávy do všech sousedních vrcholů.
V další iteraci je pak vybrán zase jiný vrchol, který bude posílat zprávy.
Iterace končí ve chvíli, kdy se už nemění marginální pravděpodobnosti proměnných.

Vzhledem k tomu, že posílané zprávy už neodpovídají skutečným faktorům, je důležitá otázka, zda-li Loopy Belief Propagation vůbec nalezne správné marginální pravděpodobnosti.
Existují podmínky za kterých algoritmus bude konvergovat~\cite{tatikonda2002loopy}.
V obecném případě ovšem může algoritmus konvergovat ke špatným pravděpodobnostem, anebo nemusí konvergovat vůbec a pak dochází k oscilacím.
Pokud algoritmus nekonverguje ke správným pravděpodobnostem, tak není ani možné se k nim přiblížit nebo je odvodit z oscilací~\cite{murphy1999loopy}.
V praxi ovšem k problémům většinou nedochází a inference konverguje v rozumném čase.

\begin{algorithm}
\caption{Loopy Belief Propagation}
\label{alg:lbp}
\begin{algorithmic}
\Function{LBP}{$F$, $S$}
\State{$F$ -- faktor graf}
\State{$S$ -- strategie}
\State

\State{$\textsc{Init-Factor-Graph}(F)$}
\Repeat
    \For{vrchol $v \in F$ vybraný podle strategie $S$}
        \For{soused $n$ vrcholu $v$}
            \If{$v$ je faktor}
                \State{$\mu_{v \rightarrow n}(n) \gets \textsc{Message-To-Var}(v, n)$}
            \Else
                \State{$\mu_{v \rightarrow n}(v) \gets \textsc{Message-To-Factor}(v, n)$}
            \EndIf
        \EndFor
    \EndFor
\Until{konvergence}

\For{každá proměnná $X$}
    \State{$P(X) = \prod_{f \in ne(X)} \mu_{f \rightarrow X}(X)$}
\EndFor
\EndFunction

\State

\Function{Init-Factor-Graph}{$F$}
\State{$F$ -- faktor graf}
\State

\For{každou proměnnou $X \in F$}
    \For{sousední faktor $f$}
        \State{$\mu_{X \rightarrow f}(X) \gets 1$}
        \State{$\mu_{f \rightarrow X}(X) \gets 1$}
    \EndFor
\EndFor
\EndFunction

\State

\Function{Message-To-Var}{$f$, $v$}
\State{$f$ -- zdrojový faktor}
\State{$v$ -- cílová proměnná}
\State

\State{$\vec{X} \gets \text{Dom}(f)$}
\State{$\mu_{f \rightarrow v}(v) \gets \sum_{\vec{X} \backslash v} f(\vec{X}) \prod_{u \in ne(f) \backslash v} \mu_{u \rightarrow f}(u)$}
\EndFunction

\State

\Function{Message-To-Factor}{$v$, $f$}
\State{$v$ -- zdrojová proměnná}
\State{$f$ -- cílový faktor}
\State

\State{$\mu_{v \rightarrow f}(v) \gets \prod_{g \in ne(v) \backslash f} \mu_{g \rightarrow v}(v)$ }
\EndFunction
\end{algorithmic}
\end{algorithm}

\subsubsection{Popis LBP}

Popis Loopy Belief Propagation je v algoritmu~\ref{alg:lbp}.
Při inferenci ve faktor grafu je nejprve třeba inicializovat všechny zprávy.
Posílání zpráv budeme iterovat dokud zprávy nedokonvergují, anebo můžeme nastavit pevný počet iterací.
Pak vybíráme vrcholy a z každého pošleme zprávu do všech okolních.
Způsobem výběru vrcholů se budeme zabývat později.
I náhodný výběr vrcholů povede k výsledku, dodání externí informace či chytřejší heuristiky ovšem může výpočet výrazně urychlit.

Poslání zprávy se liší podle toho, zda-li posíláme z proměnné do faktoru, anebo naopak.
V obou případech vynásobíme příchozí zprávy ze všech sousedních vrcholů, kromě toho, do kterého zprávu posíláme.
Při posílání zprávy z vecholu pro faktor zprávy vynásobíme faktorem a marginalizujeme všechny proměnné, kromě té, které zprávu posíláme.

Ve faktor grafu se tedy zprávy šíří pouze přes pravděpodobnostní rozdělení jedné proměnné.
To právě může vést k oscilacím anebo konvergenci ke špatným hodnotám. 
Příkladem je například skrytý markovský model se dvěma skrytými binárními proměnnými a pozorováním pro každou skrytou proměnnou.
Pokud je mezi proměnnými závislost XOR, pak nikdy nebudeme pozorovat $(1, 1)$, ale přesto bude každá proměnná konvergovat k rovnoměrnému rozdělení. A tedy sdružená pravděpodobnost pro pozorování $(1, 1)$ dostane nenulovou pravděpodobnost.

Při počítání zpráv často musíme pronásobit spoustu zpráv z okolních vrcholů.
Přitom se zprávy zas tak často nemění.
Častou optimalizací je pro každý vrchol si pamatovat součin všech jeho zpráv
\begin{equation}
    b(v) = \begin{cases}
        \prod_{u \in ne(v)} \mu_{u \rightarrow v}(v) & \mbox{$v$ je proměnná} \\
        f_v(\vec{X}) \prod_{u \in ne(v)} \mu_{u \rightarrow v}(u) & \mbox{$v$ je faktor} \\
    \end{cases}
\end{equation}
Pak můžeme při odesílání zprávy z vrcholu $v$ do vrcholu $u$ pouze $b(v)$ vydělit zprávou $\mu_{u \rightarrow v}$, pro faktory marginalizovat, a zpráva je připravena k odeslání.
$b(v)$ můžeme aktualizovat po výběru vrcholu $v$ pro posílání zpráv, anebo po každém přijetí zprávy.

Pro pozorované proměnné je třeba metodu \textsc{Message-To-Factor} upravit tak, aby pouze posílala pozorovanou hodnotu.

Implementace Loopy Belief Algoritmu i se strategiemi pro výběr vrcholů bude popsána v kapitole~\ref{ch:kap4}.

\section{Propagace s aproximovanými zprávami}

V předchozí sekci byla faktorizována sdružená pravděpodobnost a tak došlo k zjednodušení inference.
Mezi jednotlivými vrcholy byly stále posílány exaktní zprávy.
V této sekci se zaměříme na inferenci v modelech, kde není možné spočítat zprávy exaktně a je třeba je aproximovat.
Tento problém nastává u některých modelů se spojitými náhodnými proměnnými.

Předpokládáme, že máme pravděpodobnostní grafický model, reprezentující sdruženou pravděpodobnost dat a pozorování pomocí součinu faktorů.
\begin{equation}
P(\vec{X}, \vec{E}) = \prod_i f_i(\vec{X})
\end{equation}
Nás zajímá aposteriorní distribuce $P(\vec{X} \mid \vec{E})$ pro zjištění stavu, stejně tak jako pravděpodobnost pozorování $P(\vec{E})$ pro normalizaci.
Aposteriorní distribuci vyjádříme ze sdružené.
\begin{equation}
P(\vec{X} \mid \vec{E}) = \frac{1}{P(\vec{E})} \prod_i f_i(\vec{X})
\end{equation}
a pravděpodobnost pozorování je dána
\begin{equation}
P(\vec{E}) = \int \prod_i f_i(\vec{X}) \mathrm{d}\vec{X}
\end{equation}

Pro aproximativní inferenci ukážeme algoritmus Expectation propagation (EP)~\cite{minka2001expectation}.
Vytvoříme aproximaci aposteriorní distribuce, která je také dána součinem faktorů
\begin{equation}
q(\vec{X}) = \frac{1}{Z} \prod_i \tilde{f}_i(\vec{X}),
\end{equation}
kde každý faktor $\tilde{f}_i$ je aproximace odpovídající skutečnému faktoru $f_i$ a faktor $\frac{1}{Z}$ je normalizační konstanta.
Aproximované faktory musíme zkonstruovat tak, abychom byli schopni provádět inferenci.

Pro měření vzdálenosti aproximovaného rozdělení od skutečného používáme Kullback-Leiblerovu divergenci (KL)~\cite{kullback1997information}.
KL divergence, také známá jako relativní entropie, mezi dvěma pravděpodobnostními rozděleními $p(x)$ a $q(x)$ je
\begin{equation}
KL(p \| q) = \int p(x) \log \left(\frac{p(x)}{q(x)}\right) \mathrm{d}x.
\end{equation}
Divergence splňuje tři vlastnosti:
\begin{enumerate}
\item $KL(p \| p) = 0$,
\item $KL(p \| q) = 0$ právě tehdy když $p = q$,
\item $KL(p \| q) > 0$ pro všechna $p, q$.
\end{enumerate}

\subsection{Rodina exponenciálních rozdělení}
\label{sec:expfam}

Rodina exponenciálních rozdělení~\cite{bernardo2009bayesian} přes $\vec{X}$, dáno parametry $\vec\eta$ je definována jako množina distribucí ve tvaru
\begin{equation}
P(\vec{X} \mid \vec\eta) = h(\vec{X}) g(\vec\eta) \exp \big(\vec\eta^T u(\vec{X})\big),
\label{eq:gj}
\end{equation}
kde $\vec{X}$ jsou náhodné proměnné, které mohou být diskrétní nebo spojité, $\vec\eta$ jsou nazývány přirozené parametry rozdělení, $u(\vec{X})$ je funkce $\vec{X}$.
Funkce $g(\vec\eta)$ může být interpretována jako koeficient, který zajišťuje, že je distribuce normalizována a tedy splňuje
\begin{equation}
g(\vec\eta) \int h(\vec{X}) \exp \big(\vec\eta^T u(\vec{X})\big) \mathrm{d}\vec{X} = 1
\label{eq:g1}
\end{equation}

Výhody distribucí z exponenciální rodiny si ukážeme na KL divergenci.
Nechť počítáme divergenci $KL(p \| q)$, kde $p(\vec{X})$ je zafixovaná distribuce a $q(\vec{X})$ je distribuce z exponenciální rodiny.
Pokud zapíšeme KL divergenci jako funkci $\vec\eta$, tak dostaneme
\begin{equation}
KL(p \| q) = - \log g(\vec\eta) - \vec\eta^T \mathbb{E}_{p(\vec{X})}[u(\vec{X})] + \mathrm{const}.
\end{equation}
KL divergenci můžeme minimalizovat tak, že nastavíme první derivaci podle $\vec\eta$ rovnou nule, z čehož dostaneme
\begin{equation}
-\nabla \log g(\vec\eta) = \mathbb{E}_{p(\vec{X})}[u(\vec{X})].
\label{eq:kl}
\end{equation}

Nyní se musíme podívat na $- \nabla \log g(\vec\eta)$.
Vezmeme derivaci obou stran (\ref{eq:g1}) podle $\vec\eta$.
Získáme
\begin{align}
& \nabla g(\vec\eta) \int h(\vec{X}) \exp \big(\vec\eta^T u(\vec{X}) \big) \mathrm{d}\vec{X} \\
& + g(\vec\eta) \int h(\vec{X}) \exp \big(\vec\eta^T u(\vec{X})\big) u(\vec{X}) \mathrm{d}\vec{X} = 0 \nonumber
\end{align}
Přeuspořádáním a znova použitím (\ref{eq:gj}) a (\ref{eq:g1}) získáme
\begin{equation}
-\frac{1}{g(\vec\eta)}\nabla g(\vec\eta)
= g(\vec\eta) \int h(\vec{X}) \exp\big(\vec\eta^T u(\vec{X})\big) u(\vec{X}) \mathrm{d}\vec{X} = \mathbb{E}[u(\vec{X})].
\end{equation}
Ve výsledku získáváme 
\begin{equation}
- \nabla \log g(\vec\eta) = \mathbb{E}[u(\vec{X})]
\label {eq:ge}
\end{equation}

Nový poznatek o funkci $g(\vec\eta)$ z (\ref{eq:ge}) můžeme substituovat do (\ref{eq:kl}) a získáme
\begin{equation}
\mathbb{E}_{q(\vec{X})}[u(\vec{X})] = \mathbb{E}_{p(\vec{X})}[u(\vec{X})],
\label{eq:ee}
\end{equation}
z čehož vidíme, že pro minimalizaci KL divergence nám stačí najít parametry rozdělení $q(\vec{X})$ tak, aby mělo stejné očekávané statistiky jako $p(\vec{X})$.

\subsection{Expectation Propagation}
\label{sec:ep}

Expectation propagation je algoritmus založený na aproximaci aposteriorní distribuce součinem aproximovaných faktorů
\begin{equation}
q(\vec{X}) = \frac{1}{Z} \prod_i \tilde{f}_i(\vec{X}).
\end{equation}
Z praktických důvodů uvedených v předchozí sekci jsou aproximované faktory z exponenciální rodiny.
Díky tomu bude jejich součin také z exponenciální rodiny.

Aproximující distribuci bychom chtěli nalézt pomocí minimalizace KL divergence mezi skutečnou a aproximovanou distribucí.
\begin{equation}
KL(p \| q) = KL\bigg(\frac{1}{p(\vec{E})} \prod_i f_i(\vec{X}) \bigg\| \frac{1}{Z} \prod_i \tilde{f}_i(\vec{X} \bigg)
\end{equation}
Pro počítání KL divergence ovšem potřebujeme umět počítat se skutečnou pravděpodobností a když už s ní umíme efektivně počítat, proč ztrácet čas s aproximací.

Další možností je aproximovat jednotlivé faktory, díky tomu bychom v jednom kroku nalezli aproximace pro všechny faktory a měli bychom hotovo.
Naším cílem je ale nalézt nejlepší aproximaci celé aposteriorní pravděpodobnosti a to není v tomto případě zaručeno.

Expectation propagation sice aproximuje jednotlivé faktory, ale vždy v kontextu všech ostatních~\cite{bishop2006pattern}.
Nejprve jsou všechny inicializovány a pak jsou procházeny jeden po druhém a každý je aktualizován.
Tento přístup je podobný metodám, které byly ukázány v předešlých sekcí.
Předpokládejme, že chceme aktualizovat faktor $\tilde{f}_j (\vec{X})$.
Nejprve odstraníme faktor z produktu $\prod_{i \ne j} \tilde{f}_i(\vec{X})$.
Následně nalezneme novou hodnotu pro faktor $\tilde{f}_j$ tak, aby pravděpodobnostní rozložení
\begin{equation}
q^{new}(\vec{X}) \propto \tilde{f}_j(\vec{X}) \prod_{i \ne j} \tilde{f}_i(\vec{X})
\label{eq:qnew}
\end{equation}
bylo co nejblíže 
\begin{equation}
f_j(\vec{X})\prod_{i \ne j} \tilde{f}_i(\vec{X}).
\end{equation}
Všechny faktory $i \ne j$ necháváme zafixované. 
Díky tomu zajistíme, že nová aproximace faktoru je nejpřesnější v oblastech, které mají největší aposteriorní pravděpodobnost definovanou zbývajícími faktory.

Odstranění faktoru $\tilde{f}_j(\vec{X})$ z aktuální aproximace aposteriorní distribuce provedeme vytvořením nenormalizované distribuce
\begin{equation}
q^{\backslash j}(\vec{X}) = \frac{q(\vec{X})}{\tilde{f}_j(\vec{X})}
\end{equation}
Mohli bychom počítat $q^{\backslash j}(\vec{X})$ jako součin všech faktorů kromě $j$, ale v praxi je dělení rychlejší.
Této nenormalizované distribuci se říká cavity distribuce.
Její kombinací se skutečným faktorem $f_j(\vec{X})$ dostaneme distribuci
\begin{equation}
    \frac{1}{Z_j} f_j(\vec{X}) q^{\backslash j}(\vec{X}),
\label{eq:aprox}
\end{equation}
kde $Z_j$ je normalizační konstanta, dána
\begin{equation}
    Z_j = \int f_j(\vec{X}) q^{\backslash j}(\vec{X}) \mathrm{d}\vec{X}.
\label{eq:zj}
\end{equation}

Nyní můžeme nalézt novou hodnotu pro faktor $\tilde{f}_j$ minimalizací KL divergence
\begin{equation}
    KL\bigg(\frac{f_j(\vec{X}) q^{\backslash j}(\vec{X})}{Z_j} \bigg\| q^{new}(\vec{X})\bigg).
\end{equation}
To uděláme jednoduše použitím poznatku z (\ref{eq:ee}), který říká, že nám stačí najít parametry $q^{new}$ tak, aby postačující statistiky odpovídali momentům aproximovaného rozdělení (\ref{eq:aprox}).

Z nalezené aproximace $q^{new}$ pak získáme novou hodnotu faktoru $\tilde{f}_j$ z (\ref{eq:qnew}) vydělením zbývajících faktorů
\begin{equation}
    \tilde{f}_j(\vec{X}) = K \frac{q^{new}(\vec{X})}{q^{\backslash j}(\vec{X})}
\label{eq:newfj}
\end{equation}

Koeficient $K$ získáme tak, že obě strany (\ref{eq:newfj}) vynásobíme $q^{\backslash j}(\vec{X})$ a zintegrováním, čímž získáme
\begin{equation}
    K = \int \tilde{f}_j(\vec{X}) q^{\backslash j}(\vec{X}) \mathrm{d}\vec{X},
\end{equation}
použili jsme navíc toho, že $q^{new}(\vec{X})$ je normalizovaná distribuce.
Hodnota $K$ pak může být nalezena srovnáním
\begin{equation}
    \int \tilde{f}_j (\vec{X}) q^{\backslash j}(\vec{X}) \mathrm{d}\vec{X} = \int f_j(\vec{X}) q^{\backslash j}(\vec{X}) \mathrm{d}\vec{X}
\end{equation}
Z čehož zjistíme, že $K = Z_j$ a tedy může být nalezeno přímo z (\ref{eq:zj}).

Expectation Propagation je popsán v algoritmu~\ref{alg:ep}.
Dokud aproximované rozdělení nekonverguje, vybíráme jednotlivé faktory a aktualizujeme je podle popsaného postupu.

\begin{algorithm}
\caption{Expectation propagation}
\label{alg:ep}
\begin{algorithmic}
\State Inicializujeme všechny aproximativní faktory $\tilde{f}_i(\vec{X})$ na neinformativní.
\State Inicializujeme aposteriorní aproximace nastavením: $$q(\vec{X}) \propto \prod_i \tilde{f}_i(\vec{X})$$
\Repeat
    \State Vyber faktor $\tilde{f}_j(\vec{X})$, který bude aktualizován.
    \State Odeber $\tilde{f}_j(\vec{X})$ z aposteriorní distribuce vydělením: $$q^{\backslash j} (\vec{X}) = \frac{q(\vec{X})}{\tilde{f}_j(\vec{X})}$$
    \State Vypočítej novou aposteriorní distribuci nastavením momentů $q^{new}$ na 
    \State hodnotu momentů $q^{\backslash j}(\vec{X})f_j(\vec{X})$, včetně vyhodnocení normalizační 
    \State konstanty $$Z_j = \int q^{\backslash j}(\vec{X}) f_j(\vec{X}) \mathrm{d}\vec{X}$$
    \State Nastav novou hodnotu faktoru: $$\tilde{f}_j(\vec{X}) = Z_j \frac{q^{new}(\vec{X})}{q^{\backslash j}(\vec{X})}$$
\Until{konvergence}
\State Spočítej aproximaci pozorování: $$p(\vec{E}) \simeq \int \prod_i \tilde{f}_i(\vec{X}) \mathrm{d}\vec{X}$$
\end{algorithmic}
\end{algorithm}

V praxi je třeba provést několik iterací, a v každé aktualizovat všechny faktory.
Stejně jako u LBP zde nemáme žádnou garanci, že bude algoritmus konvergovat.
Pro aproximace $q(\vec{X})$ v exponenciální rodině, pokud iterace konverguje, pak nalezené řešení je stacionární bod specifické potenciální energie~\cite{minka2001expectation}.

% !TEX root = thesis.tex
\chapter{Učení parametrů}
\label{ch:ep}

Představená metoda LBP v~minulé kapitole funguje pro modely, kde máme nastavené parametry faktorů.
V~této kapitole si představíme grafický model pro diskrétní proměnné, u~kterého jsou i parametry faktorů proměnné a je možné pro ně použít Bayesovský přístup.
Díky tomu bude možné nalézt aposteriorní distribuci pro tyto parametry a naučit je podle dat.

Inferenci v~tomto upraveném modelu už není možné dělat pomocí LBP a~budeme muset použít algoritmus Expectation Propagation.

\section{Grafický model}

Model se skládá z~faktorů a proměnných.
Nechť máme vybraný faktor $f_\beta$, spojený s~několika proměnnými
$\vec{x} = (x_0, x_1, \dots, x_{N_x})$
a množinami parametrů
$\vec{\Theta} = (\vec{\theta}_1, \dots, \vec{\theta}_{N_\theta})$.
Tento faktor reprezentuje podmíněnou pravděpodobnost:
$$f_\beta(\vec{x}, \Theta) = p(x_0 | x_1, \dots, x_{N_x}; \Theta).$$

Rodičovské proměnné $x_1, \dots, x_{N_x}$ označujeme jako $\vec{x^\prime}$.
Vektor $\vec{x^\prime}$ určuje, jaká množina parametrů bude použita.
Protože množiny parametrů jsou číslovány $1, \dots, N_\theta$ a rodičovské
proměnné $1, \dots, N_x$, musí být pro vybrání správné množiny parametrů
použito mapování $\rho(\vec{x^\prime})$.
Faktor pak může být zapsán zkráceně:
$$f_\beta(\vec{x}, \Theta) = p(x_0 | x_1, \dots, x_{N_x}; \Theta) =
\theta_{\rho(x^\prime), x_0}$$

Příklad faktoru $f_\beta$ je na obrázku~\ref{ex:factor}, pro ilustraci pouze s~třemi parametry a~čtyřmi proměnnými.
Správně by měl faktor mít tolik parametrů, kolik je možných přiřazení pro rodičovské proměnné.

\begin{figure}[h]
\begin{center}
\begin{tikzpicture}[->, main/.style={circle, draw}, factor/.style={draw}]
    
\node[main] at (-3.7, 2)    (t1) {$\vec\theta_1$};
\node[main] at (-4, 0)      (t2) {$\vec\theta_2$};
\node[main] at (-3.7, -2)   (t3) {$\vec\theta_3$};
\node[factor] at (0, 0)     (f)  {$f_\beta$};
\node[main] at (3.7, 2)     (x1) {$X_1$};
\node[main] at (4, 0)       (x2) {$X_2$};
\node[main] at (3.7, -2)    (x3) {$X_3$};
\node[main] at (0, -4)      (x0) {$X_0$};

\path
	(t1) edge [above right, near start] node {$q^{\backslash \beta}(\vec\theta_1)$} (f)
	(t2) edge [above, near start] node {$q^{\backslash \beta}(\vec\theta_2)$} (f)
	(t3) edge [below right, near start] node {$q^{\backslash \beta}(\vec\theta_3)$} (f)
	(x1) edge [above left, near start] node {$q^{\backslash \beta}(X_1)$} (f)
	(x2) edge [above, near start] node {$q^{\backslash \beta}(X_2)$} (f)
	(x3) edge [below left, near start] node {$q^{\backslash \beta}(X_3)$} (f)
	(f) edge [right, near end] node {$m_{f_\beta \rightarrow X_0}(X_0)$} (x0);


\end{tikzpicture}
\end{center}
\caption{Vybraný faktor $f_\beta$ pro aktualizaci při učení parametrů.}
\label{ex:factor}
\end{figure}

\section{Výpočet marginálních pravděpodobností}

Pro výpočet sdružené pravděpodobnosti používáme plně faktorizovanou distribuci.
Pro každou proměnnou (množinu parametrů) je její marginální pravděpodobnost
rovna součinu zpráv ze sousedních faktorů.
Pro daný faktor $f_\beta$ je neúplná distribuce $q^{\backslash \beta}(x_i)$, popř.
$q^{\backslash \beta}(\vec{\theta}_i)$, rovna součinu zpráv ze všech ostatních faktorů do $x_i$, popř $\vec\theta_i$.
Aproximovaná marginální pravděpodobnost proměnné (množiny parametrů) je pak součinem neúplné
distribuce a zprávy z~faktoru:
$$q(x_i) = q^{\backslash \beta}(x_i) m_{f_\beta \rightarrow x_i}(x_i), 
\quad \quad
q(\vec{\theta}_i) = q^{\backslash \beta}(\vec{\theta}_i) m_{f_\beta \rightarrow \theta}(\vec{\theta}_i)$$

Tyto neúplné distribuce jsou právě zprávy z~proměnné, popř. množiny parametrů, do faktoru:

$$m_{x_i \rightarrow f_\beta} = q^{\backslash \beta}(x_i), 
\quad \quad
m_{\vec\theta_i \rightarrow f_\beta} = q^{\backslash \beta}(\vec\theta_i).$$

\subsection{Marginální pravděpodobnost proměnných}

Pokud chceme aktualizovat hodnotu naší aproximace marginální pravděpodobnosti,
je třeba minimalizovat její vzdálenost od skutečné marginální
pravděpodobnosti, ve tvaru
\begin{align}
p^*(\tilde{x}_j) &=
\sum_{\vec{x} \backslash x_j} \,
	\int_{\vec{\Theta}} \,
    		\prod_i \,
			q^{\backslash \beta}(x_i) \,
		\prod_l \,
			q^{\backslash \beta}(\vec{\theta}_l) \;
		f_\beta(\vec{x};\,
    		  \vec{\Theta})
\label{eq:1} 
\\
&=
\sum_{\vec{x} \backslash x_j} \,
	\prod_i \,
		q^{\backslash \beta}(x_i) \,
    \int_{\vec{\theta_{\rho(\vec{x^\prime})}}} \,
	    q^{\backslash \beta}(\vec{\theta}_{\rho(\vec{x^\prime})})\;
    \theta_{\rho(\vec{x^\prime}), x_0} \label{eq:2} 
\\
&= 
\sum_{\vec{x} \backslash x_j} \,
	\prod_i \,
		q^{\backslash \beta}(x_i)\;
    		\mathbb{E}_{q^{\backslash \beta}} 
			(\theta_{\rho(\vec{x^\prime}), x_0}).
\label{eq:3}
\end{align}

Rovnost~(\ref{eq:1}) vychází z~definice výpočtu marginální pravděpodobnosti ze
sdružené pravděpodobnosti.
V~(\ref{eq:2}) byla použita definice faktoru, z~integrálu byly vytknuty členy neobsahující $\Theta$ a nakonec bylo využito toho, že pro
množiny parametrů, které nejsou spojeny s~faktorem $f_\beta$, je jejich neúplná
distribuce rovna marginální distribuci, a tedy $\int_{\theta_i} q(\theta_i) =
1$. V~(\ref{eq:3}) byla použita definice očekávané hodnoty.

Marginální pravděpodobnost proměnné $x_j$ tedy je
\begin{equation}
p^*(x_j) =
\sum_{\vec{x} \backslash x_j} \;
	\prod_i \;
		q^{\backslash \beta}(x_i)\;
    		\mathbb{E}_{q^{\backslash \beta}} 
			(\theta_{\rho(\vec{x^\prime}), x_0}).
\end{equation}

Tady docházíme k~výsledku, který je velmi podobný výpočtu marginální
pravděpodobnosti v~Loopy Belief Propagation algoritmu.

Zprávu z~faktoru $f_\beta$ do vrcholu $x_j$ pak získáme vydělením zprávy z~$x_j$ do $f_\beta$ z~marginální pravděpodobnosti,

\begin{equation}
m_{f_\beta \rightarrow x_j}(x_j) =
    \sum_{\vec{x} \backslash x_j} \;
        \prod_{i \ne j}\;
            q^{\backslash \beta}(x_i)\;
            \mathbb{E}_{q^{\backslash \beta}}
                (\theta_{\rho(\vec{x^\prime}), x_0}).
\label{eq:msgfromftox}
\end{equation}

\subsection{Marginální pravděpodobnost parametrů}

Pro množiny parametrů se jejich marginální pravděpodobnost spočítá podobně jako
pro proměnné:
\begin{align}
p^*(\tilde{\vec{\theta}}_j) & = 
    \sum_{\vec{x}} \; 
    \int_{\vec{\Theta} \backslash \vec{\theta}_j} \;
         \prod_i \;
             q^{\backslash \beta}(x_i) \;
         \prod_l \;
    q^{\backslash \beta}(\vec{\theta}_l) \;
    f_\beta(\vec{x}; \vec{\Theta}) \label{eq:ep:theta_1}
\\
& = 
    \sum_{l \ne j} \;
        \sum_{\vec{x}: \rho(\vec{x^\prime}) = l} \;
            \prod_i \;
                q^{\backslash \beta}(x_i) \,
                \int_{\vec{\Theta} \backslash \vec{\theta}_j} \;
                    \prod_k \;
                        q^{\backslash \beta}(\vec{\theta}_k) \;
                        \theta_{l, x_0}\; + \label{eq:ep:theta_2}
\\
& + 
    \sum_{\vec{x}: \rho(\vec{x^\prime}) = j} \;
        \prod_i \;
            q^{\backslash \beta}(x_i) \;
            \int_{\vec{\Theta} \backslash \vec{\theta}_j} \;
                \prod_k \;
                    q^{\backslash \beta}(\vec{\theta}_k) \; 
                    \tilde{\theta}_{j, x_0}
\nonumber
\\
& = 
    \left[ \;
        \sum_{l \ne j} \;
            \sum_{\vec{x}: \rho(\vec{x^\prime}) = l} \;
                \prod_i \;
                    q^{\backslash \beta}(x_i) \;
                    \mathbb{E}_{q^{\backslash \beta}(\vec{\theta}_l)} (\theta_{l, x_0}) \;
    \right] \;
    q^{\backslash \beta}(\tilde{\vec{\theta}}_j) \; + \label{eq:ep:theta_3}
\\
& + 
    \sum_{\vec{x}: \rho(\vec{x^\prime}) = j} \;
        \prod_i \;
            q^{\backslash \beta}(x_i) \;
            \tilde{\theta}_{j,x_0} \;
            q^{\backslash \beta}(\tilde{\vec{\theta}}_j) \;
\label{eq:ep:theta_4}
\nonumber
\\
& = w_0 \; q^{\backslash \beta}(\tilde{\vec{\theta}}_j) \, + \, \sum_k \; w_k \;
    \tilde{\theta}_{j,k} \; q^{\backslash \beta}(\tilde{\vec{\theta}}_j),
\end{align}

kde

\begin{align}
w_0 &=
	\sum_{l \ne j} \;
		\sum_{\vec{x}: \rho(\vec{x^\prime}) = l} \;
			\prod_i \;
 				q^{\backslash \beta}(x_i) \;
				\mathbb{E}_{q^{\backslash \beta}(\vec{\theta}_l)} (
					\theta_{l,
					x_0}) 
\\
w_k &=
	\sum_{\vec{x}: \rho(\vec{x^\prime}) = j, x_0 = k} \; 
		\prod_i \;
			q^{\backslash \beta}(x_i)		
\end{align}

Opět vycházíme z~výpočtu marginální pravděpodobnosti podle definice.
V~rovnici \eqref{eq:ep:theta_2} jsme rozdělili sumu přes
$\vec{x}$ na ty, pro které se ve faktoru použije množina parametrů
$\tilde{\vec{\theta}}_j$, a na ty ostatní. Také jsme z~integrálu vytknuli součin
neúplných distribucí pro proměnné. V~dalším kroku \eqref{eq:ep:theta_3} jsme opět
použili toho, že integrál přes $\vec{\Theta}$ je ve skutečnosti několik
integrálů přes jednotlivé množiny parametrů. Díky tomu je můžeme vložit mezi členy produktu neúplných distribucí pro množiny parametrů. Ve výsledku
získáme $q^{\backslash \beta}(\tilde{\vec{\theta}}_j) \, \int_{\vec{\theta}_l} \,
q^{\backslash \beta}(\vec{\theta}_l) \; \theta_{l, x_0}$ a pak zbylé členy, které zmizí.

Docházíme k~vyjádření skutečné marginální pravděpodobnosti, v~níž není třeba
integrovat přes všechny množiny parametrů, ale stačí jen očekávaná hodnota
těchto parametrů.

\section{Aproximace marginálních pravděpodobností}

Při výpočtu aproximující distribuce $q(\vec{\theta}_j)$ se může stát, že skutečná marginální distribuce je v~komplikovaném tvaru a pak je aproximace složitá na výpočet a může se odchylovat od skutečné distribuce.
Například se jedná o~směs distribucí.
Ideální případ je takový, kdy aproximace je stejného typu jako skutečná distribuce.

Pro zjednodušení výpočtů, jak se dále ukáže, zvolíme zprávy z~faktoru do množiny parametrů, $m_{f \rightarrow \vec{\theta}_i}(\vec{\theta}_i)$, ve tvaru Dirichletovské distribuce s~parametry $\vec\alpha_{f \rightarrow \vec\theta_i}$:
\begin{equation}
m_{f \rightarrow \vec{\theta}_i}(\vec{\theta}_i) =
	Dir(\vec{\theta}_i; \vec\alpha_{f \rightarrow \vec\theta_i}) =
        \frac{\Gamma (\sum_j \vec\alpha_{f \rightarrow \vec\theta_i, j})}
             {\prod_j \Gamma(\vec\alpha_{f \rightarrow \vec\theta_i, j})}
        \prod_j \theta_{i,j}^{\vec\alpha_{f \rightarrow \vec\theta_i, j} - 1},
\end{equation}
kde $\Gamma$ je Gamma funkce (zobecnění faktoriálu):
\begin{equation}
    \Gamma(z) = \int_0^\infty \! t^{z-1} \exp(-t) \;\mathrm{d}t.
\end{equation}

Dirichletovská distribuce byla zvolena, protože má důležité vlastnosti pro
součin, které budou využity dále pro výpočet neúplné distribuce a celkové
aproximace. Pokud označíme aproximované faktory indexem $\beta$ a každý bude
mít vlastní parametry $\vec\alpha_{f_\beta \rightarrow \vec\theta_i}$, tak výsledná aproximace bude
tvaru:
\begin{align}
q(\vec{\theta}_i) \,& \, \propto
    \prod_\beta \;
    m_{f_\beta \rightarrow \vec{\theta}_i}(\vec{\theta}_i) \,
 \propto \,
    \prod_\beta \;
        \prod_j \;
            \theta_{i,j}^{\vec\alpha_{f_\beta \rightarrow \vec\theta_i, j} \,-\, 1}
\\
&\propto \,
    Dir(\vec{\theta}_i ;\, \sum_\beta \vec\alpha_{f_\beta \rightarrow \vec\theta_i} - (\,|\beta| - 1\,) \, \vec{1})
\\
&=
    Dir(\vec{\theta}_i;\,
        \vec{\alpha}_i),
\label{eq:ep:aprx_1}
\end{align}
kde $\vec{\alpha}_i = \sum_\beta \vec\alpha_{f_\beta \rightarrow \vec\theta_i} - (\,|\beta| - 1\,)
\vec{1})$.

Při aktualizaci faktoru $\tilde\beta$ tedy neúplná distribuce bude
\begin{align}
q^{\backslash \tilde\beta} (\vec{\theta_i})
&\propto
    \prod_{\beta \ne \tilde\beta}
        m_{f_\beta \rightarrow \vec{\theta}_i}(\vec{\theta}_i)
\\
& \propto
Dir(\vec{\theta}_i;
    \vec{\alpha}_i - \vec\alpha_{f_\beta \rightarrow \vec\theta_i} + \vec{1}).
\end{align}

Naším cílem je nalézt parametry $\vec{\alpha}^*$ aproximované marginální pravděpodobnosti \eqref{eq:ep:aprx_1}, které minimalizují vzdálenost od skutečné marginální pravděpodobnosti \eqref{eq:ep:theta_4}. 

Jednotlivé faktory musíme vždy aproximovat s~ohledem na všechny ostatní.
U~našeho grafického modelu využijeme nezávislostí, a tak nám stačí pouze pracovat se zprávami ze sousedních vrcholů.
Všechny informace se v~modelu šíří přes zprávy mezi vrcholy, a proto je třeba několika iterací pro přesun informace z~jedné části modelu do druhé.

Pro měření vzdálenosti mezi dvěma pravděpodobnostními rozloženími se používá Kullback-Leiblerova divergence:
\begin{equation}
KL(p \| q) =
\int_{-\infty}^{\infty}
    p(x) \log\left(\frac{p(x)}{q(x)}\right) \mathrm{d}x
\end{equation}
Pro nalezení minima použijeme algoritmus srovnání momentů a budeme
minimalizovat $KL(p^*\| q)$.

Pokud se podíváme na skutečnou marginální pravděpodobnost
$p^*(\vec{\theta}_i)$, zjistíme, že můžeme některé její členy upravit.
Využijeme také vlastnosti Gamma funkce $\Gamma(x) = (x-1) \Gamma(x-1)$.

\begin{align}
w_j \, \theta_j \, Dir(\vec{\theta}; \vec{\alpha}) &\propto \,
    w_j \, \theta_j \, \frac{\Gamma (\, \sum_i \, \alpha_{i} \,)}{\prod_i \, \Gamma(\,\alpha_i\,)}\;
    \prod_i \, \theta_{i}^{\alpha_{i} - 1}
\label{eq:4}
\\
&\propto w_j \, \frac{\Gamma (\,\sum_i \, \alpha_{i}\,)} {\prod_i \,
    \Gamma(\,\alpha_{i}\,)} \; \theta_j^{\alpha_j} \prod_{i \ne j} \,
    \theta_{i}^{\alpha_{i} - 1} \\
&\propto w_j \,
    \frac{\Gamma (\,\sum_i \alpha_{i}\,)}
         {\prod_i \, \Gamma(\, \alpha_{i} \,)} \;
    \frac{\Gamma(\, \alpha_{j} + 1 \,) \; \prod_{i \ne j} \, \Gamma(\,\alpha_i\,)}
         {\Gamma (\,1 + \sum_i \alpha_i\,)} \;
    Dir(\vec\theta; \, \vec{\alpha} + \vec{\delta}_j) \\
&\propto w_j \,
    \frac{\Gamma (\,\sum_i \alpha_{i}\,)}
         {\prod_i \, \Gamma(\, \alpha_{i}\,)} \;
    \frac{\alpha_j \, \Gamma(\,\alpha_j\,) \, \prod_{i \ne j} \, \Gamma(\,\alpha_i\,)}
         {(\,\sum_i \, \alpha_i\,) \, \Gamma (\,\sum_i \,\alpha_i\,)}\;
    Dir(\vec\theta; \,\vec{\alpha} + \vec{\delta}_j) \\
&\propto w_j \,
    \frac{\alpha_j}
         {\sum_i \, \alpha_i} \;
    Dir(\vec\theta; \,\vec{\alpha} + \vec{\delta}_j) \\
\end{align}

Zde jsme použili $\vec{\delta}_i$, což je vektor nul, který má pouze na pozici $i$ jedničku.

Díky této úpravě lze $p^*$ vyjádřit jako směs Dirichletovských rozdělení.
\begin{equation}
p^*(\vec{\theta}) =
    w_0^* \; Dir(\vec{\theta};\, \vec{\alpha}) +
    \sum_j \, w^*_j \;
        Dir(\vec{\theta};\, \vec{\alpha} + \vec{\delta}_j),
\label{eq:marginaltheta}
\end{equation}
kde
\begin{align}
    w^*_0 \propto w_0, \quad \quad
    w^*_j \propto w_j \, \frac{\alpha_j}{\sum_i \alpha_i}, \quad \quad
    \sum_{i=0}^k w_i^* = 1.
\end{align}

Jak bylo ukázáno v~sekci~\ref{sec:expfam}, pro minimalizaci KL divergence mezi dvěma rozděleními z~exponenciální rodiny
stačí, pokud se budou rovnat jejich postačující statistiky. Dokážeme jednoduše
spočítat první dva momenty Dirichletovského rozdělení, a proto použijeme
aproximaci a budeme počítat pouze s~nimi a zbylé momenty zanedbáme.
Je tedy třeba nalézt střední hodnotu a rozptyl proměnných z~$p^*(\vec\theta)$:

\begin{align}
\mathbb{E}_{p^*}[\vec\theta] &= \int \vec\theta \, p^*(\vec\theta) ~ \mathrm{d}\vec\theta
\\
&= \int \vec\theta \, (
	w_0^* \, Dir(\vec{\theta};\, \vec{\alpha}) +
	\sum_j \, w^*_j \,
        	Dir(\vec{\theta};\, \vec{\alpha} + \vec{\delta}_j)
    ) ~ \mathrm{d}\vec\theta
\\
&= w_0^* \int \vec\theta \, Dir(\vec{\theta};\, \vec{\alpha}) ~ d\vec\theta +
	\sum_j w_j^* \int \vec\theta \, Dir(\vec{\theta};\, \vec{\alpha} + \vec{\delta}_j)
    ~ \mathrm{d}\vec\theta
\\
&= w_0^* \, \mathbb{E}_{Dir(\vec\alpha)}[\vec\theta] +
	\sum_j \, w_j^* \, \mathbb{E}_{Dir(\vec\alpha + \vec{\delta}_j)}[\vec\theta]
\label{eq:epweimix}
\end{align}
Střední hodnotu proměnných $\vec\theta$ podle rozdělení $p^*$ lze podle (\ref{eq:epweimix}) spočítat
jako vážený součet středních hodnot $\vec\theta$ podle jednotlivých
Dirichletovských distribucí, z~kterých se $p^*$ skládá.

První moment tedy máme spočítaný, pro výpočet rozptylu můžeme využít
definici:
\begin{equation}
	Var[\vec{\theta}] = \mathbb{E}[\vec{\theta}^2] - \mathbb{E}[\vec\theta]^2.
\end{equation}

Chybí nám tedy ještě výpočet střední hodnoty druhé mocniny proměnné
$\vec\theta$ podle $p^*$. Můžeme ji vyjádřit z~definice střední hodnoty:

\begin{align}
\mathbb{E}_{p^*}[\vec\theta^2] &= \int \vec\theta^2 \, p^*(\vec\theta) ~ d\vec\theta
\\
&= w_0^* \int \vec\theta^2 \, Dir(\vec{\theta};\, \vec{\alpha}) ~ d\vec\theta +
	\sum_j w_j^* \int \vec\theta^2 \, Dir(\vec{\theta};\, \vec{\alpha} + \vec{\delta}_j)
	 ~ d\vec\theta
\\
&= w_0^* \, \mathbb{E}_{Dir(\vec\alpha)}[\vec\theta^2] +
	\sum_j w_j^* \, \mathbb{E}_{Dir(\vec\alpha + \vec{\delta}_j)}[\vec\theta^2].
\label{eq:epweimix2}
\end{align}

Opět získáváme vážený součet středních hodnot podle Dirichletovských rozdělení.
Střední hodnotu druhé mocniny proměnné podle Dirichletovského rozdělení
lze opět jednoduše odvodit z~definice:
\begin{align}
\mathbb{E}_{Dir(\vec\alpha)}[x_i^2] &=
	\int x_i^2 \, Dir(\vec x; \, \vec\alpha) \; \mathrm{d}\vec x
\\
&=
	\int x_i^2 \, \frac{\Gamma(\,\alpha_0\,)}
			   {\prod_{j=1}^N \, \Gamma(\,\alpha_j\,)} \;
		\prod_{j=1}^N x_j^{\alpha_j - 1} \; \mathrm{d}\vec x.
\end{align}

Nyní jsme v~podobné situaci jako v~(\ref{eq:4}). Budeme postupovat stejně,
pokusíme se výraz přepsat do formy Dirichletovské distribuce nad $x_i$:

\begin{align}
\mathbb{E}_{Dir(\vec\alpha)}[x_i^2] &=
\int \frac{\Gamma(\,\alpha_0 + 2\,) \; \alpha_i \; (\,\alpha_i + 1\,)}
		{\alpha_0 \; (\,\alpha_0 + 1\,) \; \Gamma(\,\alpha_i + 2\,) \; \prod_{j \ne i} \; \Gamma(\,\alpha_j\,)} \;
    x_i^{\alpha_i + 1} \, \prod_{j \ne i} \, x_j^{\alpha_j - 1} \;
\mathrm{d}\vec x
\\
&= \frac{\alpha_i \; (\,\alpha_i + 1\,)}
		{\alpha_0 \; (\,\alpha_0 + 1\,)} \,
	\int \frac{\Gamma(\,\beta_0\,)}
		{\prod_i \,\Gamma(\,\beta_i\,)} \;
        \prod_i \, x_i^{\beta_i - 1} \;
    \mathrm{d}\vec x
\\
&= \frac{\alpha_i \; (\,\alpha_i + 1\,)}
		{\alpha_0 \; (\,\alpha_0 + 1\,)}
	\int Dir(\vec x;\, \vec\beta)\;
    \mathrm{d}\vec x
\\
&= \frac{\alpha_i \; (\,\alpha_i + 1\,)}
		{\alpha_0 \; (\,\alpha_0 + 1\,)}.
\end{align}

Vyjádřili jsme $\Gamma(\,\alpha_0\,)$ a $\Gamma(\,\alpha_i\,)$ s~pomocí $\Gamma(\,\alpha_0 + 2\,)$ a $\Gamma(\,\alpha_i + 2\,)$:
\begin{align}
\Gamma(\alpha_0) = \frac{\Gamma(\,\alpha_0 + 2\,)}
				{\alpha_0\;(\,\alpha_0 + 1\,)},
\quad \quad
\Gamma(\alpha_i) = \frac{\Gamma(\,\alpha_i + 2\,)}
				{\alpha_i \; (\,\alpha_i + 1\,)}.
\end{align}

Následně jsme vytvořili nové parametry $\vec\beta$:
\begin{align}
\beta_i = \alpha_i + 2, \quad \quad
\beta_{j \ne i} = \alpha_j, \quad \quad
\beta_0 = \sum_i \beta_i.
\end{align}

Nyní tedy dokážeme spočítat $\mathbb{E}_{p^*}[\vec\theta]$ a $\mathbb{E}_{p^*}[\vec\theta^2]$.
Parametry aproximovaného rozdělení nalezneme následovně
\begin{align}
\frac{\mathbb{E}[X_1] - \mathbb{E}[X_1^2]}
     {\mathbb{E}[X_1^2] - \mathbb{E}[X_1]^2} &=
\frac{
	\frac{\alpha_1}
		{\alpha_0}
	- \frac{\alpha_1(\alpha_1 + 1)}
		{\alpha_0(\alpha_0 + 1)}
	}
	{
	\frac{\alpha_1(\alpha_1 + 1)}
		{\alpha_0(\alpha_0 + 1)}
	- \frac{\alpha_1^2}
		{\alpha_0^2}
	}
\label{eq:5}
\\
&=
\frac{
    \frac{\alpha_1 (\alpha_0 + 1) - \alpha_1(\alpha_1 + 1)}
         {\alpha_0 (\alpha_0 + 1)}
}{
    \frac{\alpha_0 \alpha_1 (\alpha_1 + 1) - \alpha_1^2 (\alpha_0 + 1)}
         {\alpha_0^2 (\alpha_0 + 1)}
}
\\
&=
\frac{\alpha_0 \alpha_1 (\alpha_0 - \alpha_1)}
     {\alpha_1 (\alpha_0 \alpha_1 + \alpha_0 - \alpha_0 \alpha_1 - \alpha_1)}
\\
&=
\alpha_0
\\
\alpha_i &= \mathbb{E}[X_i] \;\alpha_0 \label{eq:6}
\end{align}

Z~rovnice~(\ref{eq:5}) vypočítáme sumu všech parametrů $\alpha_0$. 
Protože střední hodnota proměnné z~Dirichletovského rozdělení je právě
$\frac{\alpha_i}{\alpha_0}$, tak jednotlivé parametry získáme z~rovnice~(\ref{eq:6}).

\section{Algoritmus}

V~algoritmu~\ref{alg:ep} je předveden algoritmus Expectation Propagation upravený pro potřeby učení parametrů.

\begin{algorithm}
\caption{Expectation Propagation pro učení parametrů}
\label{alg:ep}
\begin{algorithmic}
\State Parametry zpráv z~faktoru $\beta$ do množiny parametrů
    $\vec{\theta}_i$ označíme $\vec\alpha_{f_\beta \rightarrow \vec\theta_i}$.
\State Parametry zpráv z~množiny parametrů $\vec{\theta}_i$ do faktoru $\beta$
    označíme $\vec\alpha_{ \vec\theta_i \rightarrow f_\beta}$.
\State Parametry marginální distribuce množiny parametrů $\vec{\theta}_i$
    označíme $\vec{\alpha}_i$.
\Init
    \State Nastav zprávy mezi faktory a proměnnými na $1$.
    \State Nastav parametry $\vec\alpha_{f_\beta \rightarrow \vec\theta_i}$ na $1$.
    \State Nastav parametry $\vec{\alpha}_i$ na apriorní hodnotu.
\EndInit
\Repeat
    \State Vyber faktor $f_{\tilde\beta}$, který se bude aktualizovat.

    \State Spočítej všechny zprávy z~parametrů:
    \For{každý parametr $\vec{\theta}_i$ spojený s~faktorem $f_{\tilde\beta}$}
        \State Parametry zprávy z~$\vec{\theta}_i$ do $f_{\tilde\beta}$:
            $\vec\alpha_{ \vec\theta_i \rightarrow f_{\tilde\beta}} = \vec{\alpha}_i -
            \vec\alpha_{f_{\tilde\beta} \rightarrow \vec\theta_i } + 1$.
    \EndFor

    \State Aktualizuj zprávy z~faktoru do proměnných:
    \For{každou proměnnou $X_i$, spojenou s~faktorem $f_{\tilde\beta}$}
    \State Zpráva z~$f_{\tilde\beta}$ do $X_i$ podle (\ref{eq:msgfromftox}):
    \State
        \hskip\algorithmicindent
        $m_{f_{\tilde{\beta}} \rightarrow x_j} =
            \sum_{\vec{x} \backslash x_j}
                \mathbb{E}_{q^{\backslash\tilde\beta}}(\,\theta_{\rho(\vec{x^\prime}), x_0}\,)\,
                \prod_{i \ne j}\,
                    q^{\backslash\tilde\beta}(x_i)$
    \EndFor

    \State Aktualizuj marginální pravděpodobnost parametrů:
    \For{každý parametr $\vec{\theta}_i$, spojený s~faktorem $f_{\tilde\beta}$}
        \State Spočítej parametry $\vec{\alpha}_i^*$ pro Dirichletovské
            rozdělení, které nejlépe
        \State aproximuje cílovou marginální distribuci
            (\ref{eq:marginaltheta}). Metoda popsána
        \State v~předchozí sekci.
        \State Parametry zprávy z~$f_{\tilde\beta}$ do $\vec{\theta}_i$:
        \State
            \hskip\algorithmicindent
            $\vec\alpha_{f_{\tilde\beta} \rightarrow \vec\theta_i} =
                \vec{\alpha}_i^*
                - \vec\alpha_{ \vec\theta_i \rightarrow f_{\tilde\beta}}
                + 1$
        \State Aktualizuj parametry marginální distribuce $q(\vec{\theta}_i)$
        \State
            \hskip\algorithmicindent
            $\vec{\alpha}_i = \vec{\alpha}_i^* =
                \vec\alpha_{f_{\tilde\beta} \rightarrow \vec\theta_i}
                + \vec\alpha_{ \vec\theta_i \rightarrow f_{\tilde\beta}}$
    \EndFor
    \For{každou proměnnou $X_i$, spojenou s~faktorem $f_{\tilde\beta}$}
        \State Aktualizuj zprávy z~proměnných do faktoru:
        \State
            \hskip\algorithmicindent
            $q^{\backslash\tilde\beta}(x_i) =
                \prod_{\beta \ne \tilde\beta} \, m_{f_{\beta} \rightarrow x_i}(x_i)$
    \EndFor
\Until{konvergence}
\end{algorithmic}
\end{algorithm}


% Ukázka použití některých konstrukcí LateXu (odkomentujte, chcete-li)
% \include{ukazka}

% !TEX root = prace.tex
\chapter*{Závěr}
\addcontentsline{toc}{chapter}{Závěr}

Práce splnila všechny vytyčené cíle.
Výsledkem je implementace metod pro odhad stavu a parametrů v~dialogových systémech.
Byly představeny metody pro inferenci v~grafických modelech založené na posílání zpráv, navíc byl odvozen algoritmus pro inferenci parametrů dirichletovského rozdělení srovnáním momentů.
Práce obsahuje příklady použití vytvořené knihovny a otestování na reálných datech v~rámci soutěže Dialog State Tracking Challenge, kde generativní systém pro odhad stavu používající implementovanou knihovnu dosáhl výsledků srovnatelných s~ostatními přihlášenými týmy, které se dají považovat za state of the art.
Bližší detaily podle cílů vytyčených v~úvodu práce:
\begin{enumerate}

\item V~kapitole \ref{ch:kap1} byly představeny dialogové systémy a jejich jednotlivé součásti.
V~sekci \ref{sec:bn} pak byly představeny bayesovské sítě, které se nabízí jako ideální model pro reprezentaci dialogového stavu.

\item V~kapitole \ref{ch:kap2} jsou prezentovány možné inferenční algoritmy pro odhad stavu v~bayesovských sítích.
V~sekci \ref{sec:lbp} je představen algoritmus Loopy Belief Propagation.
Jeho použití je ukázáno v~sekci \ref{sec:usage}.

\item Implementace LBP byla použita pro DSTC a popis celé soutěže i s~výsledky je v~sekci \ref{sec:dstc}.
Implementovaný systém pro odhad stavu byl lepší než baseline systém poskytnutý organizátory (tabulka \ref{t:all:datasets}) a systém se navíc umístil příznivě i vzhledem k~ostatním přihlášeným týmům (tabulka \ref{t:DSTC:ranking}).

\item Implementované strategie pro výběr vrcholu v~LBP algoritmu jsou popsány v~sekci \ref{sec:noch}.

\item Algoritmus Expectation Propagation pro inferenci v~obecných grafických modelech je představen v~sekci \ref{sec:ep}.

\item Knihovna a jednotlivé její součásti jsou popsány v~kapitole \ref{ch:kap4}.
Příklad vrcholů vytvořených speciálně pro použití v~EP je v~sekci \ref{sec:vrdir}.
Učení parametrů pro odhad stavu v~dialogových systémech je odvozeno v~kapitole \ref{ch:ep}.
\end{enumerate}

\singlespacing

%%% Seznam použité literatury
\addcontentsline{toc}{chapter}{Seznam použité literatury}
\bibliography{citace}{}
\bibliographystyle{czech}

%%% Tabulky v diplomové práci, existují-li.
\chapwithtoc{Seznam tabulek}

%%% Použité zkratky v diplomové práci, existují-li, včetně jejich vysvětlení.
\chapwithtoc{Seznam použitých zkratek}

%%% Přílohy k diplomové práci, existují-li (různé dodatky jako výpisy programů,
%%% diagramy apod.). Každá příloha musí být alespoň jednou odkazována z vlastního
%%% textu práce. Přílohy se číslují.
\chapwithtoc{Přílohy}

\openright
\end{document}
